\section{核验锚点簇：天启、崇祯：辽东、灾荒与政权危机（1621--1644）}
\label{register:1621-1644_tianqi_chongzhen}
辽东防线、灾荒地区、邮递与起事网络相互牵连，却各有不同材料密度。崇祯朝没有存世实录；因此每一条都把后出汇编、奏疏和地方记忆的界限留在文字里。

以下为读者可跳转的首次描写。它们是按王朝年序编排的证据性短条，而非把账本字段伪装成同一种长篇叙事：每项只在所列材料可承受的范围内给出行动者、空间、后果与限度。

\begin{description}
  \item[\eventanchor{ML-1621-001}{1621（天启元年）五月4日：沉辽之战：后金夺沉阳}] 明、后金（清）与辽东诸势力；朝廷与官员、军队与卫所、边地政权与社群；边地接触区；不把族名或部名当固定疆界。「沉辽之战：后金夺沉阳」将朝廷与官员、军队与卫所、边地政权与社群的选择落在边地接触区；编年记录可显示先后，却不能代替对兵源、道路和行政准备的证明，故不把相邻事件直接写成原因。 「沉辽之战：后金夺沉阳」使明、后金（清）与辽东诸势力在边地接触区的下一步选择重新围绕朝廷与官员、军队与卫所、边地政权与社群展开；其影响应落到随后可见的军令、安置或交涉，而非以未经互校的人数概括。 材料：《熹宗实录》天启元年条；《明史·兵志》及相关列传；边镇、地方志或对方材料。限度：就「沉辽之战：后金夺沉阳」而论，《熹宗实录》天启元年条；《明史·兵志》及相关列传；边镇、地方志或对方材料提供的是1621（天启元年）五月4日的记录线索；朝廷与官员、军队与卫所、边地政权与社群在边地接触区的实际行动细节不能由这一材料链独自补足。
  \item[\eventanchor{ML-1621-002}{1621（天启元年）fall：社安叛乱：四川、贵州彝族叛乱}] 明与地方反抗、流动作战集团；朝廷与官员、军队与卫所、地方社会与精英、流动与反抗集团；地方或省域；未具城镇者不补造路线。发生在地方或省域的「社安叛乱：四川、贵州彝族叛乱」牵动朝廷与官员、军队与卫所、地方社会与精英、流动与反抗集团；已知材料让我们看见这一时点，却没有把准备行动、地方资源与决策意图完整连起来，因果只能保留为待核问题。 「社安叛乱：四川、贵州彝族叛乱」并未结束朝廷与官员、军队与卫所、地方社会与精英、流动与反抗集团在地方或省域的压力，而是把问题转到后续防务、运输或行政处置；只有后续记录能说明它实际造成何种改变。 材料：《熹宗实录》天启元年条；《明史·兵志》及相关列传；边镇、地方志或对方材料。限度：《熹宗实录》天启元年条；《明史·兵志》及相关列传；边镇、地方志或对方材料为「社安叛乱：四川、贵州彝族叛乱」留下编年入口，却未必保留地方或省域中朝廷与官员、军队与卫所、地方社会与精英、流动与反抗集团的完整视角；本条因此不据之裁定人数、责任或动机。
  \item[\eventanchor{ML-1621-003}{1621（天启元年）十二月：镇江堡之战：明军对后金的袭击被击退}] 明、后金（清）与辽东诸势力；朝廷与官员、军队与卫所、边地政权与社群；边地接触区；不把族名或部名当固定疆界。「镇江堡之战：明军对后金的袭击被击退」出现时，朝廷与官员、军队与卫所、边地政权与社群是在边地接触区的限制下行动；资料所给的是年序定位而非动机自述，前期军政压力只能作为解释线索，不能充作定论。 在边地接触区，朝廷与官员、军队与卫所、边地政权与社群因「镇江堡之战：明军对后金的袭击被击退」而要重新安排守备、行军或沟通；这一后续压力可被标示，但战果与伤亡不能只凭一方叙事定量。 材料：《熹宗实录》天启元年条；《明史·兵志》及相关列传；边镇、地方志或对方材料。限度：「镇江堡之战：明军对后金的袭击被击退」借《熹宗实录》天启元年条；《明史·兵志》及相关列传；边镇、地方志或对方材料获得可检索的时间坐标；边地接触区里朝廷与官员、军队与卫所、边地政权与社群的路线、人数与后果若无其他材料相证，均保留为不可见的部分。
  \item[\eventanchor{ML-1621-004}{1621（天启元年）：万历末至天启初的立储、内阁和言路材料标记中枢运作的堵塞与调整，派系归类需回到具体文书。}] 明；朝廷与官员；跨区域行政、资源或军政网络；不假定同步执行。继承、官僚协调或皇权运作的压力使问题进入朝廷程序。 它影响后续人事与政策议程；动机和派系标签不能由单一叙事裁定。 材料：《熹宗实录》天启元年条；《明史·本纪》及相关列传；诏令、奏疏或会典版本。限度：桥接条目只标示可回查的年度问题与材料链；实录有中央选择性，崇祯期尤其需要奏疏、地方、朝鲜及满文材料交叉。
  \item[\eventanchor{ML-1621-005}{1621（天启元年）：本年灾害、疫病或地方赈济线索应以受灾地区文书和地方志复核，中央奏报不能覆盖全部社会。}] 明；朝廷与官员、地方社会与精英；地方或省域；未具城镇者不补造路线。自然风险与地方报告促成朝廷或地方的应对记录。 赈济、迁徙和损失的实际范围仍须以受灾地区材料分别核验。 材料：《熹宗实录》天启元年条；《明史·五行志》；受灾府县地方志与奏疏。限度：桥接条目只标示可回查的年度问题与材料链；实录有中央选择性，崇祯期尤其需要奏疏、地方、朝鲜及满文材料交叉。
  \item[\eventanchor{ML-1622-001}{1622（天启二年）三月11日：广宁之战：后金夺取广宁}] 明、后金（清）与辽东诸势力；朝廷与官员、军队与卫所、边地政权与社群；边地接触区；不把族名或部名当固定疆界。发生在边地接触区的「广宁之战：后金夺取广宁」牵动朝廷与官员、军队与卫所、边地政权与社群；已知材料让我们看见这一时点，却没有把准备行动、地方资源与决策意图完整连起来，因果只能保留为待核问题。 作为后续影响，「广宁之战：后金夺取广宁」改变了朝廷与官员、军队与卫所、边地政权与社群在边地接触区继续调度、驻守或交涉的起点；能确认的是条件变化，而不是由单一叙述推算出的完整战果。 材料：《熹宗实录》天启二年条；《明史·兵志》及相关列传；边镇、地方志或对方材料。限度：「广宁之战：后金夺取广宁」可由《熹宗实录》天启二年条；《明史·兵志》及相关列传；边镇、地方志或对方材料定位到1622（天启二年）三月11日，但这些叙述未完整保存边地接触区的执行与受影响者经验，不能把沉默写成事实。
  \item[\eventanchor{ML-1622-002}{1622（天启二年）六月24日：澳门之战：荷兰人对澳门的进攻被葡萄牙人击退}] 明与海上、朝贡相关政权；朝廷与官员、军队与卫所；账本可辨空间；地点细节仍须回到所列材料。朝贡名义、边境安全与港口贸易的利益使交涉持续发生。 它为后续册封、互市或海防安排留下文书与跨政权比较线索。 材料：《熹宗实录》天启二年条；《明史·外国传》；《朝鲜王朝实录》、琉球《历代宝案》或海外档案。限度：「澳门之战：荷兰人对澳门的进攻被葡萄牙人击退」在1622（天启二年）六月24日的出现可回查《熹宗实录》天启二年条；《明史·外国传》；《朝鲜王朝实录》、琉球《历代宝案》或海外档案；它们不等于账本可辨空间的全景记录，朝廷与官员、军队与卫所未被记下的选择与损失不作推定。
  \item[\eventanchor{ML-1622-003}{1622（天启二年）六月：山东出现白莲教叛乱分子}] 明与地方反抗、流动作战集团；朝廷与官员、军队与卫所、地方社会与精英、流动与反抗集团、宗教与文化行动者；地方或省域；未具城镇者不补造路线。继承、官僚协调或皇权运作的压力使问题进入朝廷程序。 它影响后续人事与政策议程；动机和派系标签不能由单一叙事裁定。 材料：《熹宗实录》天启二年条；《明史·本纪》及相关列传；诏令、奏疏或会典版本。限度：《熹宗实录》天启二年条；《明史·本纪》及相关列传；诏令、奏疏或会典版本只能把「山东出现白莲教叛乱分子」置入1622（天启二年）六月的年序；对于地方或省域中朝廷与官员、军队与卫所、地方社会与精英、流动与反抗集团、宗教与文化行动者的路线、人数和动机，仍须另据地方或对方材料核验。
  \item[\eventanchor{ML-1622-004}{1622（天启二年）十月：中荷冲突：荷兰船只开始袭击明商船}] 明与海上、朝贡相关政权；朝廷与官员、财政与商业行动者；账本可辨空间；地点细节仍须回到所列材料。朝贡名义、边境安全与港口贸易的利益使交涉持续发生。 它为后续册封、互市或海防安排留下文书与跨政权比较线索。 材料：《熹宗实录》天启二年条；《明史·外国传》；《朝鲜王朝实录》、琉球《历代宝案》或海外档案。限度：「中荷冲突：荷兰船只开始袭击明商船」借《熹宗实录》天启二年条；《明史·外国传》；《朝鲜王朝实录》、琉球《历代宝案》或海外档案获得可检索的时间坐标；账本可辨空间里朝廷与官员、财政与商业行动者的路线、人数与后果若无其他材料相证，均保留为不可见的部分。
  \item[\eventanchor{ML-1622-006}{1622（天启二年）：甘肃地震已造成12000人死亡}] 明；朝廷与官员；跨区域行政、资源或军政网络；不假定同步执行。自然风险与地方报告促成朝廷或地方的应对记录。 赈济、迁徙和损失的实际范围仍须以受灾地区材料分别核验。 材料：《熹宗实录》天启二年条；《明史·五行志》；受灾府县地方志与奏疏。限度：《熹宗实录》天启二年条；《明史·五行志》；受灾府县地方志与奏疏记录「甘肃地震已造成12000人死亡」的方式限定了可说范围：可追踪1622（天启二年）的行动节点，不能凭此还原跨区域行政、资源或军政网络中朝廷与官员的每一处部署。
  \item[\eventanchor{ML-1623-001}{1623（天启三年）十月：中荷冲突：荷兰袭击厦门被击退}] 明与海上、朝贡相关政权；朝廷与官员；账本可辨空间；地点细节仍须回到所列材料。朝贡名义、边境安全与港口贸易的利益使交涉持续发生。 它为后续册封、互市或海防安排留下文书与跨政权比较线索。 材料：《熹宗实录》天启三年条；《明史·外国传》；《朝鲜王朝实录》、琉球《历代宝案》或海外档案。限度：「中荷冲突：荷兰袭击厦门被击退」借《熹宗实录》天启三年条；《明史·外国传》；《朝鲜王朝实录》、琉球《历代宝案》或海外档案获得可检索的时间坐标；账本可辨空间里朝廷与官员的路线、人数与后果若无其他材料相证，均保留为不可见的部分。
  \item[\eventanchor{ML-1623-002}{1623（天启三年）：社安之乱：明军被击败}] 明与地方反抗、流动作战集团；朝廷与官员、军队与卫所、地方社会与精英、流动与反抗集团；地方或省域；未具城镇者不补造路线。对于明与地方反抗、流动作战集团的「社安之乱：明军被击败」，地方或省域与朝廷与官员、军队与卫所、地方社会与精英、流动与反抗集团限定了行动的可行范围；史料未逐项记录前置命令和供给，不能以后来结果补造一个整齐的起因。 「社安之乱：明军被击败」改变了朝廷与官员、军队与卫所、地方社会与精英、流动与反抗集团处理地方或省域事务的顺序与代价；后续影响须由连贯记录显示，不能从孤立军报推出可量化的整体结局。 材料：《熹宗实录》天启三年条；《明史·兵志》及相关列传；边镇、地方志或对方材料。限度：「社安之乱：明军被击败」虽见于《熹宗实录》天启三年条；《明史·兵志》及相关列传；边镇、地方志或对方材料，但材料形成者未必观察到地方或省域的全部过程；朝廷与官员、军队与卫所、地方社会与精英、流动与反抗集团的行动范围和受影响群体不能由此单独决定。
  \item[\eventanchor{ML-1623-003}{1623（天启三年）：黄河决堤淹没徐州}] 明；朝廷与官员；跨区域行政、资源或军政网络；不假定同步执行。自然风险与地方报告促成朝廷或地方的应对记录。 赈济、迁徙和损失的实际范围仍须以受灾地区材料分别核验。 材料：《熹宗实录》天启三年条；《明史·五行志》；受灾府县地方志与奏疏。限度：就「黄河决堤淹没徐州」而论，《熹宗实录》天启三年条；《明史·五行志》；受灾府县地方志与奏疏提供的是1623（天启三年）的记录线索；朝廷与官员在跨区域行政、资源或军政网络的实际行动细节不能由这一材料链独自补足。
  \item[\eventanchor{ML-1623-004}{1623（天启三年）：矿税、加派、漕运和辽饷的本年奏疏与账目提供财政压力线索，预算、欠额和实解不得混同。}] 明；朝廷与官员、财政与商业行动者；跨区域行政、资源或军政网络；不假定同步执行。财政征解、交通工程或货币支付的压力使相关措施进入政务。 其法定安排不等于地方执行，后续须以地域材料追踪负担与成效。 材料：《熹宗实录》天启三年条；《明会典》相关条目；地方志、黄册/契约/河工材料。限度：桥接条目只标示可回查的年度问题与材料链；实录有中央选择性，崇祯期尤其需要奏疏、地方、朝鲜及满文材料交叉。
  \item[\eventanchor{ML-1623-005}{1623（天启三年）：万历末至天启初的立储、内阁和言路材料标记中枢运作的堵塞与调整，派系归类需回到具体文书。}] 明；朝廷与官员；跨区域行政、资源或军政网络；不假定同步执行。继承、官僚协调或皇权运作的压力使问题进入朝廷程序。 它影响后续人事与政策议程；动机和派系标签不能由单一叙事裁定。 材料：《熹宗实录》天启三年条；《明史·本纪》及相关列传；诏令、奏疏或会典版本。限度：桥接条目只标示可回查的年度问题与材料链；实录有中央选择性，崇祯期尤其需要奏疏、地方、朝鲜及满文材料交叉。
  \item[\eventanchor{ML-1624-001}{1624（天启四年）八月26日：中荷冲突：明朝军队将荷兰人逐出澎湖，他们撤退到台湾，定居在大佑安湾附近的一个海盗村附近}] 明与海上、朝贡相关政权；朝廷与官员、军队与卫所；账本可辨空间；地点细节仍须回到所列材料。朝贡名义、边境安全与港口贸易的利益使交涉持续发生。 它为后续册封、互市或海防安排留下文书与跨政权比较线索。 材料：《熹宗实录》天启四年条；《明史·外国传》；《朝鲜王朝实录》、琉球《历代宝案》或海外档案。限度：关于「中荷冲突：明朝军队将荷兰人逐出澎湖，他们撤退到台湾，定居在大佑安湾附近的一个海盗村附近」，《熹宗实录》天启四年条；《明史·外国传》；《朝鲜王朝实录》、琉球《历代宝案》或海外档案所能确证的是条目所示的时点和行动；账本可辨空间内朝廷与官员、军队与卫所的规模、意图及社会后果需要异源材料才能判断。
  \item[\eventanchor{ML-1624-002}{1624（天启四年）：社安之乱：明军击败叛军，但无法果断平息叛乱}] 明与地方反抗、流动作战集团；朝廷与官员、军队与卫所、地方社会与精英、流动与反抗集团；地方或省域；未具城镇者不补造路线。「社安之乱：明军击败叛军，但无法果断平息叛乱」所涉朝廷与官员、军队与卫所、地方社会与精英、流动与反抗集团的行动受地方或省域的空间与资源约束；年序材料保存了节点，却没有保存足以裁定出兵、守备或支援动机的全过程。 在地方或省域，朝廷与官员、军队与卫所、地方社会与精英、流动与反抗集团因「社安之乱：明军击败叛军，但无法果断平息叛乱」而要重新安排守备、行军或沟通；这一后续压力可被标示，但战果与伤亡不能只凭一方叙事定量。 材料：《熹宗实录》天启四年条；《明史·兵志》及相关列传；边镇、地方志或对方材料。限度：「社安之乱：明军击败叛军，但无法果断平息叛乱」借《熹宗实录》天启四年条；《明史·兵志》及相关列传；边镇、地方志或对方材料获得可检索的时间坐标；地方或省域里朝廷与官员、军队与卫所、地方社会与精英、流动与反抗集团的路线、人数与后果若无其他材料相证，均保留为不可见的部分。
  \item[\eventanchor{ML-1624-003}{1624（天启四年）：东林、书院、刻书和宗教活动的本年材料记录精英网络，社团存在不等于一致政治立场。}] 明；朝廷与官员、宗教与文化行动者；跨区域行政、资源或军政网络；不假定同步执行。制度、教育或知识传播的需要推动了相关行动或文本形成。 它提供制度和传播史的锚点，不能据此推断社会整体接受。 材料：《熹宗实录》天启四年条；《明史·选举志》或《艺文志》；文集、碑刻或书目材料。限度：桥接条目只标示可回查的年度问题与材料链；实录有中央选择性，崇祯期尤其需要奏疏、地方、朝鲜及满文材料交叉。
  \item[\eventanchor{ML-1624-004}{1624（天启四年）：辽东战事、边镇防务和西南兵事的本年军令记录资源调动，战报数字应与朝鲜、满文或地方材料比照。}] 明；朝廷与官员、军队与卫所、边地政权与社群、地方社会与精英；账本可辨空间；地点细节仍须回到所列材料。「辽东战事、边镇防务和西南兵事的本年军令记录资源调动，战报数字应与朝鲜、满文或地方材料比照」使朝廷与官员、军队与卫所、边地政权与社群、地方社会与精英在账本可辨空间的处境变得可见；本条未保存从征发到出动的全链条，故以可核条件而非概括性“压力”标示其背景。 对明而言，「辽东战事、边镇防务和西南兵事的本年军令记录资源调动，战报数字应与朝鲜、满文或地方材料比照」使账本可辨空间的资源和安全议题进入下一轮处置；本条所能承受的是这一转折，不是对战场损失或地方反应的完整裁断。 材料：《熹宗实录》天启四年条；《明史·兵志》及相关列传；边镇、地方志或对方材料。限度：桥接条目只标示可回查的年度问题与材料链；实录有中央选择性，崇祯期尤其需要奏疏、地方、朝鲜及满文材料交叉。
  \item[\eventanchor{ML-1624-005}{1624（天启四年）：矿税、加派、漕运和辽饷的本年奏疏与账目提供财政压力线索，预算、欠额和实解不得混同。}] 明；朝廷与官员、财政与商业行动者；跨区域行政、资源或军政网络；不假定同步执行。财政征解、交通工程或货币支付的压力使相关措施进入政务。 其法定安排不等于地方执行，后续须以地域材料追踪负担与成效。 材料：《熹宗实录》天启四年条；《明会典》相关条目；地方志、黄册/契约/河工材料。限度：桥接条目只标示可回查的年度问题与材料链；实录有中央选择性，崇祯期尤其需要奏疏、地方、朝鲜及满文材料交叉。
  \item[\eventanchor{ML-1625-001}{1625（天启五年）：东林运动被肃清}] 明；朝廷与官员；跨区域行政、资源或军政网络；不假定同步执行。继承、官僚协调或皇权运作的压力使问题进入朝廷程序。 它影响后续人事与政策议程；动机和派系标签不能由单一叙事裁定。 材料：《熹宗实录》天启五年条；《明史·本纪》及相关列传；诏令、奏疏或会典版本。限度：「东林运动被肃清」在1625（天启五年）的出现可回查《熹宗实录》天启五年条；《明史·本纪》及相关列传；诏令、奏疏或会典版本；它们不等于跨区域行政、资源或军政网络的全景记录，朝廷与官员未被记下的选择与损失不作推定。
  \item[\eventanchor{ML-1625-002}{1625（天启五年）：辽饷、剿饷、练饷及地方征解的本年文书显示财政负担，定额、征收与实际流向不得混为一谈。}] 明；朝廷与官员、军队与卫所、地方社会与精英、财政与商业行动者；地方或省域；未具城镇者不补造路线。财政征解、交通工程或货币支付的压力使相关措施进入政务。 其法定安排不等于地方执行，后续须以地域材料追踪负担与成效。 材料：《熹宗实录》天启五年条；《明会典》相关条目；地方志、黄册/契约/河工材料。限度：桥接条目只标示可回查的年度问题与材料链；实录有中央选择性，崇祯期尤其需要奏疏、地方、朝鲜及满文材料交叉。
  \item[\eventanchor{ML-1625-003}{1625（天启五年）：辽东防务、后金（清）军事行动和流动作战集团的本年记录标记多战场互动，军数和胜败须保留分歧。（材料核查重点：诏令或奏议的形成与发受链）}] 明；朝廷与官员、军队与卫所、边地政权与社群、流动与反抗集团；边地接触区；不把族名或部名当固定疆界。「辽东防务、后金（清）军事行动和流动作战集团的本年记录标记多战场互动，军数和胜败须保留分歧」出现时，朝廷与官员、军队与卫所、边地政权与社群、流动与反抗集团是在边地接触区的限制下行动；资料所给的是年序定位而非动机自述，前期军政压力只能作为解释线索，不能充作定论。 「辽东防务、后金（清）军事行动和流动作战集团的本年记录标记多战场互动，军数和胜败须保留分歧」之后，边地接触区内朝廷与官员、军队与卫所、边地政权与社群、流动与反抗集团仍须处理由该行动留下的控制、通行或防务问题；本条不把战报中的胜负、斩获或伤亡扩大为无从核对的结论。 材料：《熹宗实录》天启五年条；《明史·兵志》及相关列传；边镇、地方志或对方材料。限度：桥接条目只标示可回查的年度问题与材料链；实录有中央选择性，崇祯期尤其需要奏疏、地方、朝鲜及满文材料交叉。；本条特别核对诏令或奏议的形成与发受链。
  \item[\eventanchor{ML-1625-004}{1625（天启五年）：天启末至崇祯朝的任免、奏疏和廷议构成本年中枢危机线索；崇祯无实录，必须以多类材料互校。（材料核查重点：诏令或奏议的形成与发受链）}] 明；朝廷与官员；跨区域行政、资源或军政网络；不假定同步执行。继承、官僚协调或皇权运作的压力使问题进入朝廷程序。 它影响后续人事与政策议程；动机和派系标签不能由单一叙事裁定。 材料：《熹宗实录》天启五年条；《明史·本纪》及相关列传；诏令、奏疏或会典版本。限度：桥接条目只标示可回查的年度问题与材料链；实录有中央选择性，崇祯期尤其需要奏疏、地方、朝鲜及满文材料交叉。；本条特别核对诏令或奏议的形成与发受链。
  \item[\eventanchor{ML-1625-005}{1625（天启五年）：地方结社、宗教、出版或士人文集的本年线索提供战乱中的社会观察，幸存者记忆不能概括全国。（材料核查重点：诏令或奏议的形成与发受链）}] 明；朝廷与官员、军队与卫所、地方社会与精英、宗教与文化行动者；地方或省域；未具城镇者不补造路线。制度、教育或知识传播的需要推动了相关行动或文本形成。 它提供制度和传播史的锚点，不能据此推断社会整体接受。 材料：《熹宗实录》天启五年条；《明史·选举志》或《艺文志》；文集、碑刻或书目材料。限度：桥接条目只标示可回查的年度问题与材料链；实录有中央选择性，崇祯期尤其需要奏疏、地方、朝鲜及满文材料交叉。；本条特别核对诏令或奏议的形成与发受链。
  \item[\eventanchor{ML-1625-006}{1625（天启五年）：地方结社、宗教、出版或士人文集的本年线索提供战乱中的社会观察，幸存者记忆不能概括全国。（材料核查重点：报称信息的来源、抄传与行政分类）}] 明；朝廷与官员、军队与卫所、地方社会与精英、宗教与文化行动者；地方或省域；未具城镇者不补造路线。制度、教育或知识传播的需要推动了相关行动或文本形成。 它提供制度和传播史的锚点，不能据此推断社会整体接受。 材料：《熹宗实录》天启五年条；《明史·选举志》或《艺文志》；文集、碑刻或书目材料。限度：桥接条目只标示可回查的年度问题与材料链；实录有中央选择性，崇祯期尤其需要奏疏、地方、朝鲜及满文材料交叉。；本条特别核对报称信息的来源、抄传与行政分类。
  \item[\eventanchor{ML-1625-008}{1625（天启五年）：辽东防务、后金（清）军事行动和流动作战集团的本年记录标记多战场互动，军数和胜败须保留分歧。（材料核查重点：报称信息的来源、抄传与行政分类）}] 明；朝廷与官员、军队与卫所、边地政权与社群、流动与反抗集团；边地接触区；不把族名或部名当固定疆界。「辽东防务、后金（清）军事行动和流动作战集团的本年记录标记多战场互动，军数和胜败须保留分歧」把朝廷与官员、军队与卫所、边地政权与社群、流动与反抗集团带到边地接触区的具体关口；现存年条只留下这一行动节点，前期兵力、粮道和地方协力的次序未逐项保存，不能倒推为单一动因。 「辽东防务、后金（清）军事行动和流动作战集团的本年记录标记多战场互动，军数和胜败须保留分歧」改变了朝廷与官员、军队与卫所、边地政权与社群、流动与反抗集团处理边地接触区事务的顺序与代价；后续影响须由连贯记录显示，不能从孤立军报推出可量化的整体结局。 材料：《熹宗实录》天启五年条；《明史·兵志》及相关列传；边镇、地方志或对方材料。限度：桥接条目只标示可回查的年度问题与材料链；实录有中央选择性，崇祯期尤其需要奏疏、地方、朝鲜及满文材料交叉。；本条特别核对报称信息的来源、抄传与行政分类。
  \item[\eventanchor{ML-1626-001}{1626（天启六年）二月10日：宁远之战：后金攻宁远被击退，努尔哈赤负伤身亡}] 明、后金（清）与辽东诸势力；朝廷与官员、军队与卫所、边地政权与社群；边地接触区；不把族名或部名当固定疆界。发生在边地接触区的「宁远之战：后金攻宁远被击退，努尔哈赤负伤身亡」牵动朝廷与官员、军队与卫所、边地政权与社群；已知材料让我们看见这一时点，却没有把准备行动、地方资源与决策意图完整连起来，因果只能保留为待核问题。 「宁远之战：后金攻宁远被击退，努尔哈赤负伤身亡」并未结束朝廷与官员、军队与卫所、边地政权与社群在边地接触区的压力，而是把问题转到后续防务、运输或行政处置；只有后续记录能说明它实际造成何种改变。 材料：《熹宗实录》天启六年条；《明史·兵志》及相关列传；边镇、地方志或对方材料。限度：《熹宗实录》天启六年条；《明史·兵志》及相关列传；边镇、地方志或对方材料为「宁远之战：后金攻宁远被击退，努尔哈赤负伤身亡」留下编年入口，却未必保留边地接触区中朝廷与官员、军队与卫所、边地政权与社群的完整视角；本条因此不据之裁定人数、责任或动机。
  \item[\eventanchor{ML-1626-002}{1626（天启六年）：辽东防务、后金（清）军事行动和流动作战集团的本年记录标记多战场互动，军数和胜败须保留分歧。}] 明；朝廷与官员、军队与卫所、边地政权与社群、流动与反抗集团；边地接触区；不把族名或部名当固定疆界。「辽东防务、后金（清）军事行动和流动作战集团的本年记录标记多战场互动，军数和胜败须保留分歧」将朝廷与官员、军队与卫所、边地政权与社群、流动与反抗集团的选择落在边地接触区；编年记录可显示先后，却不能代替对兵源、道路和行政准备的证明，故不把相邻事件直接写成原因。 发生「辽东防务、后金（清）军事行动和流动作战集团的本年记录标记多战场互动，军数和胜败须保留分歧」后，边地接触区仍是朝廷与官员、军队与卫所、边地政权与社群、流动与反抗集团必须面对的行动场域；可见结果是问题被带入下一段年序，未见的伤亡、掠夺与居民去留不在本条擅断。 材料：《熹宗实录》天启六年条；《明史·兵志》及相关列传；边镇、地方志或对方材料。限度：桥接条目只标示可回查的年度问题与材料链；实录有中央选择性，崇祯期尤其需要奏疏、地方、朝鲜及满文材料交叉。
  \item[\eventanchor{ML-1626-003}{1626（天启六年）：朝鲜、辽东和海上关系的本年材料提供外部观察位置，明、后金（清）与地方势力的称谓需具体化。（材料核查重点：诏令或奏议的形成与发受链）}] 明；朝廷与官员、边地政权与社群、地方社会与精英；账本可辨空间；地点细节仍须回到所列材料。朝贡名义、边境安全与港口贸易的利益使交涉持续发生。 它为后续册封、互市或海防安排留下文书与跨政权比较线索。 材料：《熹宗实录》天启六年条；《明史·外国传》；《朝鲜王朝实录》、琉球《历代宝案》或海外档案。限度：桥接条目只标示可回查的年度问题与材料链；实录有中央选择性，崇祯期尤其需要奏疏、地方、朝鲜及满文材料交叉。；本条特别核对诏令或奏议的形成与发受链。
  \item[\eventanchor{ML-1626-004}{1626（天启六年）：辽饷、剿饷、练饷及地方征解的本年文书显示财政负担，定额、征收与实际流向不得混为一谈。（材料核查重点：诏令或奏议的形成与发受链）}] 明；朝廷与官员、军队与卫所、地方社会与精英、财政与商业行动者；地方或省域；未具城镇者不补造路线。财政征解、交通工程或货币支付的压力使相关措施进入政务。 其法定安排不等于地方执行，后续须以地域材料追踪负担与成效。 材料：《熹宗实录》天启六年条；《明会典》相关条目；地方志、黄册/契约/河工材料。限度：桥接条目只标示可回查的年度问题与材料链；实录有中央选择性，崇祯期尤其需要奏疏、地方、朝鲜及满文材料交叉。；本条特别核对诏令或奏议的形成与发受链。
  \item[\eventanchor{ML-1626-005}{1626（天启六年）：旱灾、蝗灾、疫病、河患与赈济的本年材料连接环境和社会压力，死亡与流离数字须谨慎处理。（材料核查重点：诏令或奏议的形成与发受链）}] 明；朝廷与官员；跨区域行政、资源或军政网络；不假定同步执行。自然风险与地方报告促成朝廷或地方的应对记录。 赈济、迁徙和损失的实际范围仍须以受灾地区材料分别核验。 材料：《熹宗实录》天启六年条；《明史·五行志》；受灾府县地方志与奏疏。限度：桥接条目只标示可回查的年度问题与材料链；实录有中央选择性，崇祯期尤其需要奏疏、地方、朝鲜及满文材料交叉。；本条特别核对诏令或奏议的形成与发受链。
  \item[\eventanchor{ML-1626-006}{1626（天启六年）：旱灾、蝗灾、疫病、河患与赈济的本年材料连接环境和社会压力，死亡与流离数字须谨慎处理。（材料核查重点：报称信息的来源、抄传与行政分类）}] 明；朝廷与官员；跨区域行政、资源或军政网络；不假定同步执行。自然风险与地方报告促成朝廷或地方的应对记录。 赈济、迁徙和损失的实际范围仍须以受灾地区材料分别核验。 材料：《熹宗实录》天启六年条；《明史·五行志》；受灾府县地方志与奏疏。限度：桥接条目只标示可回查的年度问题与材料链；实录有中央选择性，崇祯期尤其需要奏疏、地方、朝鲜及满文材料交叉。；本条特别核对报称信息的来源、抄传与行政分类。
  \item[\eventanchor{ML-1626-007}{1626（天启六年）：辽饷、剿饷、练饷及地方征解的本年文书显示财政负担，定额、征收与实际流向不得混为一谈。（材料核查重点：报称信息的来源、抄传与行政分类）}] 明；朝廷与官员、军队与卫所、地方社会与精英、财政与商业行动者；地方或省域；未具城镇者不补造路线。财政征解、交通工程或货币支付的压力使相关措施进入政务。 其法定安排不等于地方执行，后续须以地域材料追踪负担与成效。 材料：《熹宗实录》天启六年条；《明会典》相关条目；地方志、黄册/契约/河工材料。限度：桥接条目只标示可回查的年度问题与材料链；实录有中央选择性，崇祯期尤其需要奏疏、地方、朝鲜及满文材料交叉。；本条特别核对报称信息的来源、抄传与行政分类。
  \item[\eventanchor{ML-1626-008}{1626（天启六年）：朝鲜、辽东和海上关系的本年材料提供外部观察位置，明、后金（清）与地方势力的称谓需具体化。（材料核查重点：报称信息的来源、抄传与行政分类）}] 明；朝廷与官员、边地政权与社群、地方社会与精英；账本可辨空间；地点细节仍须回到所列材料。朝贡名义、边境安全与港口贸易的利益使交涉持续发生。 它为后续册封、互市或海防安排留下文书与跨政权比较线索。 材料：《熹宗实录》天启六年条；《明史·外国传》；《朝鲜王朝实录》、琉球《历代宝案》或海外档案。限度：桥接条目只标示可回查的年度问题与材料链；实录有中央选择性，崇祯期尤其需要奏疏、地方、朝鲜及满文材料交叉。；本条特别核对报称信息的来源、抄传与行政分类。
  \item[\eventanchor{ML-1627-001}{1627（天启七年）春：宁晋之战：皇太极率领的后金军队进攻锦州，但被击退}] 明、后金（清）与辽东诸势力；朝廷与官员、军队与卫所、边地政权与社群；边地接触区；不把族名或部名当固定疆界。「宁晋之战：皇太极率领的后金军队进攻锦州，但被击退」记录了朝廷与官员、军队与卫所、边地政权与社群在边地接触区的一次处置；其兵员、道路或地方支持并非都见于同一材料，因而不以总括理由覆盖未见环节。 「宁晋之战：皇太极率领的后金军队进攻锦州，但被击退」之后，边地接触区内朝廷与官员、军队与卫所、边地政权与社群仍须处理由该行动留下的控制、通行或防务问题；本条不把战报中的胜负、斩获或伤亡扩大为无从核对的结论。 材料：《熹宗实录》天启七年条；《明史·兵志》及相关列传；边镇、地方志或对方材料。限度：《熹宗实录》天启七年条；《明史·兵志》及相关列传；边镇、地方志或对方材料使「宁晋之战：皇太极率领的后金军队进攻锦州，但被击退」的时间位置可追查；至于朝廷与官员、军队与卫所、边地政权与社群在边地接触区的命令传递、路线和损失，必须保持待核而不能随编年补写。
  \item[\eventanchor{ML-1627-002}{1627（天启七年）九月30日：天启皇帝驾崩}] 明；朝廷与官员；跨区域行政、资源或军政网络；不假定同步执行。继承、官僚协调或皇权运作的压力使问题进入朝廷程序。 它影响后续人事与政策议程；动机和派系标签不能由单一叙事裁定。 材料：《熹宗实录》天启七年条；《明史·本纪》及相关列传；诏令、奏疏或会典版本。限度：《熹宗实录》天启七年条；《明史·本纪》及相关列传；诏令、奏疏或会典版本只能把「天启皇帝驾崩」置入1627（天启七年）九月30日的年序；对于跨区域行政、资源或军政网络中朝廷与官员的路线、人数和动机，仍须另据地方或对方材料核验。
  \item[\eventanchor{ML-1627-003}{1627（天启七年）十月2日：朱由检即位崇祯皇帝}] 明；朝廷与官员；跨区域行政、资源或军政网络；不假定同步执行。继承、官僚协调或皇权运作的压力使问题进入朝廷程序。 它影响后续人事与政策议程；动机和派系标签不能由单一叙事裁定。 材料：《熹宗实录》天启七年条；《明史·本纪》及相关列传；诏令、奏疏或会典版本。限度：「朱由检即位崇祯皇帝」在1627（天启七年）十月2日的出现可回查《熹宗实录》天启七年条；《明史·本纪》及相关列传；诏令、奏疏或会典版本；它们不等于跨区域行政、资源或军政网络的全景记录，朝廷与官员未被记下的选择与损失不作推定。
  \item[\eventanchor{ML-1627-004}{1627（天启七年）：辽东防务、后金（清）军事行动和流动作战集团的本年记录标记多战场互动，军数和胜败须保留分歧。（材料核查重点：诏令或奏议的形成与发受链）}] 明；朝廷与官员、军队与卫所、边地政权与社群、流动与反抗集团；边地接触区；不把族名或部名当固定疆界。发生在边地接触区的「辽东防务、后金（清）军事行动和流动作战集团的本年记录标记多战场互动，军数和胜败须保留分歧」牵动朝廷与官员、军队与卫所、边地政权与社群、流动与反抗集团；已知材料让我们看见这一时点，却没有把准备行动、地方资源与决策意图完整连起来，因果只能保留为待核问题。 「辽东防务、后金（清）军事行动和流动作战集团的本年记录标记多战场互动，军数和胜败须保留分歧」并未结束朝廷与官员、军队与卫所、边地政权与社群、流动与反抗集团在边地接触区的压力，而是把问题转到后续防务、运输或行政处置；只有后续记录能说明它实际造成何种改变。 材料：《熹宗实录》天启七年条；《明史·兵志》及相关列传；边镇、地方志或对方材料。限度：桥接条目只标示可回查的年度问题与材料链；实录有中央选择性，崇祯期尤其需要奏疏、地方、朝鲜及满文材料交叉。；本条特别核对诏令或奏议的形成与发受链。
  \item[\eventanchor{ML-1627-005}{1627（天启七年）：天启末至崇祯朝的任免、奏疏和廷议构成本年中枢危机线索；崇祯无实录，必须以多类材料互校。（材料核查重点：诏令或奏议的形成与发受链）}] 明；朝廷与官员；跨区域行政、资源或军政网络；不假定同步执行。继承、官僚协调或皇权运作的压力使问题进入朝廷程序。 它影响后续人事与政策议程；动机和派系标签不能由单一叙事裁定。 材料：《熹宗实录》天启七年条；《明史·本纪》及相关列传；诏令、奏疏或会典版本。限度：桥接条目只标示可回查的年度问题与材料链；实录有中央选择性，崇祯期尤其需要奏疏、地方、朝鲜及满文材料交叉。；本条特别核对诏令或奏议的形成与发受链。
  \item[\eventanchor{ML-1627-006}{1627（天启七年）：天启末至崇祯朝的任免、奏疏和廷议构成本年中枢危机线索；崇祯无实录，必须以多类材料互校。（材料核查重点：报称信息的来源、抄传与行政分类）}] 明；朝廷与官员；跨区域行政、资源或军政网络；不假定同步执行。继承、官僚协调或皇权运作的压力使问题进入朝廷程序。 它影响后续人事与政策议程；动机和派系标签不能由单一叙事裁定。 材料：《熹宗实录》天启七年条；《明史·本纪》及相关列传；诏令、奏疏或会典版本。限度：桥接条目只标示可回查的年度问题与材料链；实录有中央选择性，崇祯期尤其需要奏疏、地方、朝鲜及满文材料交叉。；本条特别核对报称信息的来源、抄传与行政分类。
  \item[\eventanchor{ML-1627-007}{1627（天启七年）：辽东防务、后金（清）军事行动和流动作战集团的本年记录标记多战场互动，军数和胜败须保留分歧。（材料核查重点：报称信息的来源、抄传与行政分类）}] 明；朝廷与官员、军队与卫所、边地政权与社群、流动与反抗集团；边地接触区；不把族名或部名当固定疆界。「辽东防务、后金（清）军事行动和流动作战集团的本年记录标记多战场互动，军数和胜败须保留分歧」将朝廷与官员、军队与卫所、边地政权与社群、流动与反抗集团的选择落在边地接触区；编年记录可显示先后，却不能代替对兵源、道路和行政准备的证明，故不把相邻事件直接写成原因。 「辽东防务、后金（清）军事行动和流动作战集团的本年记录标记多战场互动，军数和胜败须保留分歧」使明在边地接触区的下一步选择重新围绕朝廷与官员、军队与卫所、边地政权与社群、流动与反抗集团展开；其影响应落到随后可见的军令、安置或交涉，而非以未经互校的人数概括。 材料：《熹宗实录》天启七年条；《明史·兵志》及相关列传；边镇、地方志或对方材料。限度：桥接条目只标示可回查的年度问题与材料链；实录有中央选择性，崇祯期尤其需要奏疏、地方、朝鲜及满文材料交叉。；本条特别核对报称信息的来源、抄传与行政分类。
  \item[\eventanchor{ML-1627-008}{1627（天启七年）：地方结社、宗教、出版或士人文集的本年线索提供战乱中的社会观察，幸存者记忆不能概括全国。}] 明；朝廷与官员、军队与卫所、地方社会与精英、宗教与文化行动者；地方或省域；未具城镇者不补造路线。制度、教育或知识传播的需要推动了相关行动或文本形成。 它提供制度和传播史的锚点，不能据此推断社会整体接受。 材料：《熹宗实录》天启七年条；《明史·选举志》或《艺文志》；文集、碑刻或书目材料。限度：桥接条目只标示可回查的年度问题与材料链；实录有中央选择性，崇祯期尤其需要奏疏、地方、朝鲜及满文材料交叉。
  \item[\eventanchor{ML-1628-001}{1628（崇祯元年）春：山西遭遇干旱}] 明与各地社会、边疆政权；朝廷与官员；跨区域行政、资源或军政网络；不假定同步执行。自然风险与地方报告促成朝廷或地方的应对记录。 赈济、迁徙和损失的实际范围仍须以受灾地区材料分别核验。 材料：《崇祯长编》崇祯元年条及《明史·思宗本纪》；《明史·五行志》；受灾府县地方志与奏疏。限度：「山西遭遇干旱」虽见于《崇祯长编》崇祯元年条及《明史·思宗本纪》；《明史·五行志》；受灾府县地方志与奏疏，但材料形成者未必观察到跨区域行政、资源或军政网络的全部过程；朝廷与官员的行动范围和受影响群体不能由此单独决定。
  \item[\eventanchor{ML-1628-003}{1628（崇祯元年）：地方结社、宗教、出版或士人文集的本年线索提供战乱中的社会观察，幸存者记忆不能概括全国。}] 明与各地社会、边疆政权；朝廷与官员、军队与卫所、地方社会与精英、宗教与文化行动者；地方或省域；未具城镇者不补造路线。制度、教育或知识传播的需要推动了相关行动或文本形成。 它提供制度和传播史的锚点，不能据此推断社会整体接受。 材料：《崇祯长编》崇祯元年条及《明史·思宗本纪》；《明史·选举志》或《艺文志》；文集、碑刻或书目材料。限度：桥接条目只标示可回查的年度问题与材料链；实录有中央选择性，崇祯期尤其需要奏疏、地方、朝鲜及满文材料交叉。
  \item[\eventanchor{ML-1628-004}{1628（崇祯元年）：朝鲜、辽东和海上关系的本年材料提供外部观察位置，明、后金（清）与地方势力的称谓需具体化。（材料核查重点：诏令或奏议的形成与发受链）}] 明与各地社会、边疆政权；朝廷与官员、边地政权与社群、地方社会与精英；账本可辨空间；地点细节仍须回到所列材料。朝贡名义、边境安全与港口贸易的利益使交涉持续发生。 它为后续册封、互市或海防安排留下文书与跨政权比较线索。 材料：《崇祯长编》崇祯元年条及《明史·思宗本纪》；《明史·外国传》；《朝鲜王朝实录》、琉球《历代宝案》或海外档案。限度：桥接条目只标示可回查的年度问题与材料链；实录有中央选择性，崇祯期尤其需要奏疏、地方、朝鲜及满文材料交叉。；本条特别核对诏令或奏议的形成与发受链。
  \item[\eventanchor{ML-1628-005}{1628（崇祯元年）：辽饷、剿饷、练饷及地方征解的本年文书显示财政负担，定额、征收与实际流向不得混为一谈。（材料核查重点：诏令或奏议的形成与发受链）}] 明与各地社会、边疆政权；朝廷与官员、军队与卫所、地方社会与精英、财政与商业行动者；地方或省域；未具城镇者不补造路线。财政征解、交通工程或货币支付的压力使相关措施进入政务。 其法定安排不等于地方执行，后续须以地域材料追踪负担与成效。 材料：《崇祯长编》崇祯元年条及《明史·思宗本纪》；《明会典》相关条目；地方志、黄册/契约/河工材料。限度：桥接条目只标示可回查的年度问题与材料链；实录有中央选择性，崇祯期尤其需要奏疏、地方、朝鲜及满文材料交叉。；本条特别核对诏令或奏议的形成与发受链。
  \item[\eventanchor{ML-1628-006}{1628（崇祯元年）：辽饷、剿饷、练饷及地方征解的本年文书显示财政负担，定额、征收与实际流向不得混为一谈。（材料核查重点：报称信息的来源、抄传与行政分类）}] 明与各地社会、边疆政权；朝廷与官员、军队与卫所、地方社会与精英、财政与商业行动者；地方或省域；未具城镇者不补造路线。财政征解、交通工程或货币支付的压力使相关措施进入政务。 其法定安排不等于地方执行，后续须以地域材料追踪负担与成效。 材料：《崇祯长编》崇祯元年条及《明史·思宗本纪》；《明会典》相关条目；地方志、黄册/契约/河工材料。限度：桥接条目只标示可回查的年度问题与材料链；实录有中央选择性，崇祯期尤其需要奏疏、地方、朝鲜及满文材料交叉。；本条特别核对报称信息的来源、抄传与行政分类。
  \item[\eventanchor{ML-1628-007}{1628（崇祯元年）：朝鲜、辽东和海上关系的本年材料提供外部观察位置，明、后金（清）与地方势力的称谓需具体化。（材料核查重点：报称信息的来源、抄传与行政分类）}] 明与各地社会、边疆政权；朝廷与官员、边地政权与社群、地方社会与精英；账本可辨空间；地点细节仍须回到所列材料。朝贡名义、边境安全与港口贸易的利益使交涉持续发生。 它为后续册封、互市或海防安排留下文书与跨政权比较线索。 材料：《崇祯长编》崇祯元年条及《明史·思宗本纪》；《明史·外国传》；《朝鲜王朝实录》、琉球《历代宝案》或海外档案。限度：桥接条目只标示可回查的年度问题与材料链；实录有中央选择性，崇祯期尤其需要奏疏、地方、朝鲜及满文材料交叉。；本条特别核对报称信息的来源、抄传与行政分类。
  \item[\eventanchor{ML-1628-008}{1628（崇祯元年）：旱灾、蝗灾、疫病、河患与赈济的本年材料连接环境和社会压力，死亡与流离数字须谨慎处理。}] 明与各地社会、边疆政权；朝廷与官员；跨区域行政、资源或军政网络；不假定同步执行。自然风险与地方报告促成朝廷或地方的应对记录。 赈济、迁徙和损失的实际范围仍须以受灾地区材料分别核验。 材料：《崇祯长编》崇祯元年条及《明史·思宗本纪》；《明史·五行志》；受灾府县地方志与奏疏。限度：桥接条目只标示可回查的年度问题与材料链；实录有中央选择性，崇祯期尤其需要奏疏、地方、朝鲜及满文材料交叉。
  \item[\eventanchor{ML-1629-001}{1629（崇祯二年）冬：己巳事变：后金军队攻破长城，掠夺北京周边地区}] 明、后金（清）与辽东诸势力；朝廷与官员、军队与卫所、边地政权与社群；边地接触区；不把族名或部名当固定疆界。「己巳事变：后金军队攻破长城，掠夺北京周边地区」将朝廷与官员、军队与卫所、边地政权与社群的选择落在边地接触区；编年记录可显示先后，却不能代替对兵源、道路和行政准备的证明，故不把相邻事件直接写成原因。 发生「己巳事变：后金军队攻破长城，掠夺北京周边地区」后，边地接触区仍是朝廷与官员、军队与卫所、边地政权与社群必须面对的行动场域；可见结果是问题被带入下一段年序，未见的伤亡、掠夺与居民去留不在本条擅断。 材料：《崇祯长编》崇祯二年条及《明史·思宗本纪》；《明史·兵志》及相关列传；边镇、地方志或对方材料。限度：《崇祯长编》崇祯二年条及《明史·思宗本纪》；《明史·兵志》及相关列传；边镇、地方志或对方材料记录「己巳事变：后金军队攻破长城，掠夺北京周边地区」的方式限定了可说范围：可追踪1629（崇祯二年）冬的行动节点，不能凭此还原边地接触区中朝廷与官员、军队与卫所、边地政权与社群的每一处部署。
  \item[\eventanchor{ML-1629-002}{1629（崇祯二年）：崇祯皇帝削减帝国邮政经费，导致山西失业邮政工人叛乱}] 明与地方反抗、流动作战集团；朝廷与官员、军队与卫所、地方社会与精英、流动与反抗集团、财政与商业行动者；地方或省域；未具城镇者不补造路线。财政征解、交通工程或货币支付的压力使相关措施进入政务。 其法定安排不等于地方执行，后续须以地域材料追踪负担与成效。 材料：《崇祯长编》崇祯二年条及《明史·思宗本纪》；《明会典》相关条目；地方志、黄册/契约/河工材料。限度：「崇祯皇帝削减帝国邮政经费，导致山西失业邮政工人叛乱」可由《崇祯长编》崇祯二年条及《明史·思宗本纪》；《明会典》相关条目；地方志、黄册/契约/河工材料定位到1629（崇祯二年），但这些叙述未完整保存地方或省域的执行与受影响者经验，不能把沉默写成事实。
  \item[\eventanchor{ML-1629-003}{1629（崇祯二年）：She-An叛乱：叛军被击败}] 明与地方反抗、流动作战集团；朝廷与官员、军队与卫所、地方社会与精英、流动与反抗集团；地方或省域；未具城镇者不补造路线。「She-An叛乱：叛军被击败」出现时，朝廷与官员、军队与卫所、地方社会与精英、流动与反抗集团是在地方或省域的限制下行动；资料所给的是年序定位而非动机自述，前期军政压力只能作为解释线索，不能充作定论。 就后果说，「She-An叛乱：叛军被击败」重排了地方或省域中朝廷与官员、军队与卫所、地方社会与精英、流动与反抗集团可以据守、经过或谈判的条件；现有记载不能据此精确判定兵数、伤亡或地方社会的全部代价。 材料：《崇祯长编》崇祯二年条及《明史·思宗本纪》；《明史·兵志》及相关列传；边镇、地方志或对方材料。限度：以《崇祯长编》崇祯二年条及《明史·思宗本纪》；《明史·兵志》及相关列传；边镇、地方志或对方材料核对「She-An叛乱：叛军被击败」时，只能确认其在1629（崇祯二年）的叙述位置；朝廷与官员、军队与卫所、地方社会与精英、流动与反抗集团在地方或省域的具体经验仍需由可定位的异类材料补证。
  \item[\eventanchor{ML-1629-004}{1629（崇祯二年）：地方结社、宗教、出版或士人文集的本年线索提供战乱中的社会观察，幸存者记忆不能概括全国。（材料核查重点：诏令或奏议的形成与发受链）}] 明与各地社会、边疆政权；朝廷与官员、军队与卫所、地方社会与精英、宗教与文化行动者；地方或省域；未具城镇者不补造路线。制度、教育或知识传播的需要推动了相关行动或文本形成。 它提供制度和传播史的锚点，不能据此推断社会整体接受。 材料：《崇祯长编》崇祯二年条及《明史·思宗本纪》；《明史·选举志》或《艺文志》；文集、碑刻或书目材料。限度：桥接条目只标示可回查的年度问题与材料链；实录有中央选择性，崇祯期尤其需要奏疏、地方、朝鲜及满文材料交叉。；本条特别核对诏令或奏议的形成与发受链。
  \item[\eventanchor{ML-1629-005}{1629（崇祯二年）：辽东防务、后金（清）军事行动和流动作战集团的本年记录标记多战场互动，军数和胜败须保留分歧。（材料核查重点：诏令或奏议的形成与发受链）}] 明与各地社会、边疆政权；朝廷与官员、军队与卫所、边地政权与社群、流动与反抗集团；边地接触区；不把族名或部名当固定疆界。「辽东防务、后金（清）军事行动和流动作战集团的本年记录标记多战场互动，军数和胜败须保留分歧」使朝廷与官员、军队与卫所、边地政权与社群、流动与反抗集团在边地接触区的处境变得可见；本条未保存从征发到出动的全链条，故以可核条件而非概括性“压力”标示其背景。 「辽东防务、后金（清）军事行动和流动作战集团的本年记录标记多战场互动，军数和胜败须保留分歧」使明与各地社会、边疆政权在边地接触区的下一步选择重新围绕朝廷与官员、军队与卫所、边地政权与社群、流动与反抗集团展开；其影响应落到随后可见的军令、安置或交涉，而非以未经互校的人数概括。 材料：《崇祯长编》崇祯二年条及《明史·思宗本纪》；《明史·兵志》及相关列传；边镇、地方志或对方材料。限度：桥接条目只标示可回查的年度问题与材料链；实录有中央选择性，崇祯期尤其需要奏疏、地方、朝鲜及满文材料交叉。；本条特别核对诏令或奏议的形成与发受链。
  \item[\eventanchor{ML-1629-006}{1629（崇祯二年）：辽东防务、后金（清）军事行动和流动作战集团的本年记录标记多战场互动，军数和胜败须保留分歧。（材料核查重点：报称信息的来源、抄传与行政分类）}] 明与各地社会、边疆政权；朝廷与官员、军队与卫所、边地政权与社群、流动与反抗集团；边地接触区；不把族名或部名当固定疆界。从边地接触区看「辽东防务、后金（清）军事行动和流动作战集团的本年记录标记多战场互动，军数和胜败须保留分歧」，朝廷与官员、军队与卫所、边地政权与社群、流动与反抗集团的行动连接着当年的军政网络；资料没有充分说明何种先行安排直接促成它，不能把编年相邻当作因果证明。 「辽东防务、后金（清）军事行动和流动作战集团的本年记录标记多战场互动，军数和胜败须保留分歧」并未结束朝廷与官员、军队与卫所、边地政权与社群、流动与反抗集团在边地接触区的压力，而是把问题转到后续防务、运输或行政处置；只有后续记录能说明它实际造成何种改变。 材料：《崇祯长编》崇祯二年条及《明史·思宗本纪》；《明史·兵志》及相关列传；边镇、地方志或对方材料。限度：桥接条目只标示可回查的年度问题与材料链；实录有中央选择性，崇祯期尤其需要奏疏、地方、朝鲜及满文材料交叉。；本条特别核对报称信息的来源、抄传与行政分类。
  \item[\eventanchor{ML-1629-007}{1629（崇祯二年）：地方结社、宗教、出版或士人文集的本年线索提供战乱中的社会观察，幸存者记忆不能概括全国。（材料核查重点：报称信息的来源、抄传与行政分类）}] 明与各地社会、边疆政权；朝廷与官员、军队与卫所、地方社会与精英、宗教与文化行动者；地方或省域；未具城镇者不补造路线。制度、教育或知识传播的需要推动了相关行动或文本形成。 它提供制度和传播史的锚点，不能据此推断社会整体接受。 材料：《崇祯长编》崇祯二年条及《明史·思宗本纪》；《明史·选举志》或《艺文志》；文集、碑刻或书目材料。限度：桥接条目只标示可回查的年度问题与材料链；实录有中央选择性，崇祯期尤其需要奏疏、地方、朝鲜及满文材料交叉。；本条特别核对报称信息的来源、抄传与行政分类。
  \item[\eventanchor{ML-1629-008}{1629（崇祯二年）：天启末至崇祯朝的任免、奏疏和廷议构成本年中枢危机线索；崇祯无实录，必须以多类材料互校。}] 明与各地社会、边疆政权；朝廷与官员；跨区域行政、资源或军政网络；不假定同步执行。继承、官僚协调或皇权运作的压力使问题进入朝廷程序。 它影响后续人事与政策议程；动机和派系标签不能由单一叙事裁定。 材料：《崇祯长编》崇祯二年条及《明史·思宗本纪》；《明史·本纪》及相关列传；诏令、奏疏或会典版本。限度：桥接条目只标示可回查的年度问题与材料链；实录有中央选择性，崇祯期尤其需要奏疏、地方、朝鲜及满文材料交叉。
  \item[\eventanchor{ML-1630-001}{1630（崇祯三年）夏：稷巳之变：后金军队撤退}] 明、后金（清）与辽东诸势力；朝廷与官员、军队与卫所、边地政权与社群；边地接触区；不把族名或部名当固定疆界。「稷巳之变：后金军队撤退」使朝廷与官员、军队与卫所、边地政权与社群在边地接触区的处境变得可见；本条未保存从征发到出动的全链条，故以可核条件而非概括性“压力”标示其背景。 对明、后金（清）与辽东诸势力而言，「稷巳之变：后金军队撤退」使边地接触区的资源和安全议题进入下一轮处置；本条所能承受的是这一转折，不是对战场损失或地方反应的完整裁断。 材料：《崇祯长编》崇祯三年条及《明史·思宗本纪》；《明史·兵志》及相关列传；边镇、地方志或对方材料。限度：「稷巳之变：后金军队撤退」在1630（崇祯三年）夏的出现可回查《崇祯长编》崇祯三年条及《明史·思宗本纪》；《明史·兵志》及相关列传；边镇、地方志或对方材料；它们不等于边地接触区的全景记录，朝廷与官员、军队与卫所、边地政权与社群未被记下的选择与损失不作推定。
  \item[\eventanchor{ML-1630-002}{1630（崇祯三年）：朝鲜、辽东和海上关系的本年材料提供外部观察位置，明、后金（清）与地方势力的称谓需具体化。（材料核查重点：诏令或奏议的形成与发受链）}] 明与各地社会、边疆政权；朝廷与官员、边地政权与社群、地方社会与精英；账本可辨空间；地点细节仍须回到所列材料。朝贡名义、边境安全与港口贸易的利益使交涉持续发生。 它为后续册封、互市或海防安排留下文书与跨政权比较线索。 材料：《崇祯长编》崇祯三年条及《明史·思宗本纪》；《明史·外国传》；《朝鲜王朝实录》、琉球《历代宝案》或海外档案。限度：桥接条目只标示可回查的年度问题与材料链；实录有中央选择性，崇祯期尤其需要奏疏、地方、朝鲜及满文材料交叉。；本条特别核对诏令或奏议的形成与发受链。
  \item[\eventanchor{ML-1630-003}{1630（崇祯三年）：旱灾、蝗灾、疫病、河患与赈济的本年材料连接环境和社会压力，死亡与流离数字须谨慎处理。（材料核查重点：诏令或奏议的形成与发受链）}] 明与各地社会、边疆政权；朝廷与官员；跨区域行政、资源或军政网络；不假定同步执行。自然风险与地方报告促成朝廷或地方的应对记录。 赈济、迁徙和损失的实际范围仍须以受灾地区材料分别核验。 材料：《崇祯长编》崇祯三年条及《明史·思宗本纪》；《明史·五行志》；受灾府县地方志与奏疏。限度：桥接条目只标示可回查的年度问题与材料链；实录有中央选择性，崇祯期尤其需要奏疏、地方、朝鲜及满文材料交叉。；本条特别核对诏令或奏议的形成与发受链。
  \item[\eventanchor{ML-1630-005}{1630（崇祯三年）：朝鲜、辽东和海上关系的本年材料提供外部观察位置，明、后金（清）与地方势力的称谓需具体化。（材料核查重点：报称信息的来源、抄传与行政分类）}] 明与各地社会、边疆政权；朝廷与官员、边地政权与社群、地方社会与精英；账本可辨空间；地点细节仍须回到所列材料。朝贡名义、边境安全与港口贸易的利益使交涉持续发生。 它为后续册封、互市或海防安排留下文书与跨政权比较线索。 材料：《崇祯长编》崇祯三年条及《明史·思宗本纪》；《明史·外国传》；《朝鲜王朝实录》、琉球《历代宝案》或海外档案。限度：桥接条目只标示可回查的年度问题与材料链；实录有中央选择性，崇祯期尤其需要奏疏、地方、朝鲜及满文材料交叉。；本条特别核对报称信息的来源、抄传与行政分类。
  \item[\eventanchor{ML-1630-006}{1630（崇祯三年）：朝鲜、辽东和海上关系的本年材料提供外部观察位置，明、后金（清）与地方势力的称谓需具体化。（材料核查重点：同年地方材料所见的执行范围）}] 明与各地社会、边疆政权；朝廷与官员、军队与卫所、边地政权与社群、地方社会与精英；账本可辨空间；地点细节仍须回到所列材料。朝贡名义、边境安全与港口贸易的利益使交涉持续发生。 它为后续册封、互市或海防安排留下文书与跨政权比较线索。 材料：《崇祯长编》崇祯三年条及《明史·思宗本纪》；《明史·外国传》；《朝鲜王朝实录》、琉球《历代宝案》或海外档案。限度：桥接条目只标示可回查的年度问题与材料链；实录有中央选择性，崇祯期尤其需要奏疏、地方、朝鲜及满文材料交叉。；本条特别核对同年地方材料所见的执行范围。
  \item[\eventanchor{ML-1630-007}{1630（崇祯三年）：旱灾、蝗灾、疫病、河患与赈济的本年材料连接环境和社会压力，死亡与流离数字须谨慎处理。（材料核查重点：同年地方材料所见的执行范围）}] 明与各地社会、边疆政权；朝廷与官员、军队与卫所、地方社会与精英；地方或省域；未具城镇者不补造路线。自然风险与地方报告促成朝廷或地方的应对记录。 赈济、迁徙和损失的实际范围仍须以受灾地区材料分别核验。 材料：《崇祯长编》崇祯三年条及《明史·思宗本纪》；《明史·五行志》；受灾府县地方志与奏疏。限度：桥接条目只标示可回查的年度问题与材料链；实录有中央选择性，崇祯期尤其需要奏疏、地方、朝鲜及满文材料交叉。；本条特别核对同年地方材料所见的执行范围。
  \item[\eventanchor{ML-1631-001}{1631（崇祯四年）四月：叛军占领平度}] 明与地方反抗、流动作战集团；朝廷与官员、军队与卫所、地方社会与精英、流动与反抗集团；地方或省域；未具城镇者不补造路线。就「叛军占领平度」而言，明与地方反抗、流动作战集团与朝廷与官员、军队与卫所、地方社会与精英、流动与反抗集团在地方或省域的行动并非凭空发生；不过现存材料只标示节点，前序命令、运输和地方响应仍须逐件核验。 「叛军占领平度」之后，朝廷与官员、军队与卫所、地方社会与精英、流动与反抗集团在地方或省域的行动空间有所改写；这种变化须由随后文书和地方材料追踪，不能把初报数字直接当作最终后果。 材料：《崇祯长编》崇祯四年条及《明史·思宗本纪》；《明史·兵志》及相关列传；边镇、地方志或对方材料。限度：关于「叛军占领平度」，《崇祯长编》崇祯四年条及《明史·思宗本纪》；《明史·兵志》及相关列传；边镇、地方志或对方材料所能确证的是条目所示的时点和行动；地方或省域内朝廷与官员、军队与卫所、地方社会与精英、流动与反抗集团的规模、意图及社会后果需要异源材料才能判断。
  \item[\eventanchor{ML-1631-002}{1631（崇祯四年）十一月21日：大凌河之战：后金夺取大凌河}] 明、后金（清）与辽东诸势力；朝廷与官员、军队与卫所、边地政权与社群；边地接触区；不把族名或部名当固定疆界。「大凌河之战：后金夺取大凌河」出现时，朝廷与官员、军队与卫所、边地政权与社群是在边地接触区的限制下行动；资料所给的是年序定位而非动机自述，前期军政压力只能作为解释线索，不能充作定论。 在边地接触区，朝廷与官员、军队与卫所、边地政权与社群因「大凌河之战：后金夺取大凌河」而要重新安排守备、行军或沟通；这一后续压力可被标示，但战果与伤亡不能只凭一方叙事定量。 材料：《崇祯长编》崇祯四年条及《明史·思宗本纪》；《明史·兵志》及相关列传；边镇、地方志或对方材料。限度：「大凌河之战：后金夺取大凌河」借《崇祯长编》崇祯四年条及《明史·思宗本纪》；《明史·兵志》及相关列传；边镇、地方志或对方材料获得可检索的时间坐标；边地接触区里朝廷与官员、军队与卫所、边地政权与社群的路线、人数与后果若无其他材料相证，均保留为不可见的部分。
  \item[\eventanchor{ML-1631-003}{1631（崇祯四年）：天启末至崇祯朝的任免、奏疏和廷议构成本年中枢危机线索；崇祯无实录，必须以多类材料互校。（材料核查重点：诏令或奏议的形成与发受链）}] 明与各地社会、边疆政权；朝廷与官员；跨区域行政、资源或军政网络；不假定同步执行。继承、官僚协调或皇权运作的压力使问题进入朝廷程序。 它影响后续人事与政策议程；动机和派系标签不能由单一叙事裁定。 材料：《崇祯长编》崇祯四年条及《明史·思宗本纪》；《明史·本纪》及相关列传；诏令、奏疏或会典版本。限度：桥接条目只标示可回查的年度问题与材料链；实录有中央选择性，崇祯期尤其需要奏疏、地方、朝鲜及满文材料交叉。；本条特别核对诏令或奏议的形成与发受链。
  \item[\eventanchor{ML-1631-004}{1631（崇祯四年）：天启末至崇祯朝的任免、奏疏和廷议构成本年中枢危机线索；崇祯无实录，必须以多类材料互校。（材料核查重点：报称信息的来源、抄传与行政分类）}] 明与各地社会、边疆政权；朝廷与官员；跨区域行政、资源或军政网络；不假定同步执行。继承、官僚协调或皇权运作的压力使问题进入朝廷程序。 它影响后续人事与政策议程；动机和派系标签不能由单一叙事裁定。 材料：《崇祯长编》崇祯四年条及《明史·思宗本纪》；《明史·本纪》及相关列传；诏令、奏疏或会典版本。限度：桥接条目只标示可回查的年度问题与材料链；实录有中央选择性，崇祯期尤其需要奏疏、地方、朝鲜及满文材料交叉。；本条特别核对报称信息的来源、抄传与行政分类。
  \item[\eventanchor{ML-1631-005}{1631（崇祯四年）：地方结社、宗教、出版或士人文集的本年线索提供战乱中的社会观察，幸存者记忆不能概括全国。（材料核查重点：诏令或奏议的形成与发受链）}] 明与各地社会、边疆政权；朝廷与官员、军队与卫所、地方社会与精英、宗教与文化行动者；地方或省域；未具城镇者不补造路线。制度、教育或知识传播的需要推动了相关行动或文本形成。 它提供制度和传播史的锚点，不能据此推断社会整体接受。 材料：《崇祯长编》崇祯四年条及《明史·思宗本纪》；《明史·选举志》或《艺文志》；文集、碑刻或书目材料。限度：桥接条目只标示可回查的年度问题与材料链；实录有中央选择性，崇祯期尤其需要奏疏、地方、朝鲜及满文材料交叉。；本条特别核对诏令或奏议的形成与发受链。
  \item[\eventanchor{ML-1631-006}{1631（崇祯四年）：地方结社、宗教、出版或士人文集的本年线索提供战乱中的社会观察，幸存者记忆不能概括全国。（材料核查重点：报称信息的来源、抄传与行政分类）}] 明与各地社会、边疆政权；朝廷与官员、军队与卫所、地方社会与精英、宗教与文化行动者；地方或省域；未具城镇者不补造路线。制度、教育或知识传播的需要推动了相关行动或文本形成。 它提供制度和传播史的锚点，不能据此推断社会整体接受。 材料：《崇祯长编》崇祯四年条及《明史·思宗本纪》；《明史·选举志》或《艺文志》；文集、碑刻或书目材料。限度：桥接条目只标示可回查的年度问题与材料链；实录有中央选择性，崇祯期尤其需要奏疏、地方、朝鲜及满文材料交叉。；本条特别核对报称信息的来源、抄传与行政分类。
  \item[\eventanchor{ML-1631-007}{1631（崇祯四年）：天启末至崇祯朝的任免、奏疏和廷议构成本年中枢危机线索；崇祯无实录，必须以多类材料互校。（材料核查重点：同年地方材料所见的执行范围）}] 明与各地社会、边疆政权；朝廷与官员、军队与卫所、地方社会与精英；地方或省域；未具城镇者不补造路线。继承、官僚协调或皇权运作的压力使问题进入朝廷程序。 它影响后续人事与政策议程；动机和派系标签不能由单一叙事裁定。 材料：《崇祯长编》崇祯四年条及《明史·思宗本纪》；《明史·本纪》及相关列传；诏令、奏疏或会典版本。限度：桥接条目只标示可回查的年度问题与材料链；实录有中央选择性，崇祯期尤其需要奏疏、地方、朝鲜及满文材料交叉。；本条特别核对同年地方材料所见的执行范围。
  \item[\eventanchor{ML-1632-001}{1632（崇祯五年）二月22日：吴桥兵变：山东部队兵变攻克邓州}] 明与各地社会、边疆政权；朝廷与官员、军队与卫所、地方社会与精英；地方或省域；未具城镇者不补造路线。「吴桥兵变：山东部队兵变攻克邓州」所涉朝廷与官员、军队与卫所、地方社会与精英的行动受地方或省域的空间与资源约束；年序材料保存了节点，却没有保存足以裁定出兵、守备或支援动机的全过程。 「吴桥兵变：山东部队兵变攻克邓州」之后，地方或省域内朝廷与官员、军队与卫所、地方社会与精英仍须处理由该行动留下的控制、通行或防务问题；本条不把战报中的胜负、斩获或伤亡扩大为无从核对的结论。 材料：《崇祯长编》崇祯五年条及《明史·思宗本纪》；《明史·兵志》及相关列传；边镇、地方志或对方材料。限度：《崇祯长编》崇祯五年条及《明史·思宗本纪》；《明史·兵志》及相关列传；边镇、地方志或对方材料使「吴桥兵变：山东部队兵变攻克邓州」的时间位置可追查；至于朝廷与官员、军队与卫所、地方社会与精英在地方或省域的命令传递、路线和损失，必须保持待核而不能随编年补写。
  \item[\eventanchor{ML-1632-002}{1632（崇祯五年）：西班牙马尼拉与中国的贸易额每年达到200万比索}] 明与海上、朝贡相关政权；朝廷与官员、财政与商业行动者；账本可辨空间；地点细节仍须回到所列材料。朝贡名义、边境安全与港口贸易的利益使交涉持续发生。 它为后续册封、互市或海防安排留下文书与跨政权比较线索。 材料：《崇祯长编》崇祯五年条及《明史·思宗本纪》；《明史·外国传》；《朝鲜王朝实录》、琉球《历代宝案》或海外档案。限度：「西班牙马尼拉与中国的贸易额每年达到200万比索」虽见于《崇祯长编》崇祯五年条及《明史·思宗本纪》；《明史·外国传》；《朝鲜王朝实录》、琉球《历代宝案》或海外档案，但材料形成者未必观察到账本可辨空间的全部过程；朝廷与官员、财政与商业行动者的行动范围和受影响群体不能由此单独决定。
  \item[\eventanchor{ML-1632-003}{1632（崇祯五年）：明朝的防御规划者建造了一些星堡，但在中国并没有流行起来。}] 明与各地社会、边疆政权；朝廷与官员；跨区域行政、资源或军政网络；不假定同步执行。继承、官僚协调或皇权运作的压力使问题进入朝廷程序。 它影响后续人事与政策议程；动机和派系标签不能由单一叙事裁定。 材料：《崇祯长编》崇祯五年条及《明史·思宗本纪》；《明史·本纪》及相关列传；诏令、奏疏或会典版本。限度：「明朝的防御规划者建造了一些星堡，但在中国并没有流行起来」在1632（崇祯五年）的出现可回查《崇祯长编》崇祯五年条及《明史·思宗本纪》；《明史·本纪》及相关列传；诏令、奏疏或会典版本；它们不等于跨区域行政、资源或军政网络的全景记录，朝廷与官员未被记下的选择与损失不作推定。
  \item[\eventanchor{ML-1632-004}{1632（崇祯五年）：辽饷、剿饷、练饷及地方征解的本年文书显示财政负担，定额、征收与实际流向不得混为一谈。（材料核查重点：诏令或奏议的形成与发受链）}] 明与各地社会、边疆政权；朝廷与官员、军队与卫所、地方社会与精英、财政与商业行动者；地方或省域；未具城镇者不补造路线。财政征解、交通工程或货币支付的压力使相关措施进入政务。 其法定安排不等于地方执行，后续须以地域材料追踪负担与成效。 材料：《崇祯长编》崇祯五年条及《明史·思宗本纪》；《明会典》相关条目；地方志、黄册/契约/河工材料。限度：桥接条目只标示可回查的年度问题与材料链；实录有中央选择性，崇祯期尤其需要奏疏、地方、朝鲜及满文材料交叉。；本条特别核对诏令或奏议的形成与发受链。
  \item[\eventanchor{ML-1632-005}{1632（崇祯五年）：旱灾、蝗灾、疫病、河患与赈济的本年材料连接环境和社会压力，死亡与流离数字须谨慎处理。（材料核查重点：诏令或奏议的形成与发受链）}] 明与各地社会、边疆政权；朝廷与官员；跨区域行政、资源或军政网络；不假定同步执行。自然风险与地方报告促成朝廷或地方的应对记录。 赈济、迁徙和损失的实际范围仍须以受灾地区材料分别核验。 材料：《崇祯长编》崇祯五年条及《明史·思宗本纪》；《明史·五行志》；受灾府县地方志与奏疏。限度：桥接条目只标示可回查的年度问题与材料链；实录有中央选择性，崇祯期尤其需要奏疏、地方、朝鲜及满文材料交叉。；本条特别核对诏令或奏议的形成与发受链。
  \item[\eventanchor{ML-1632-006}{1632（崇祯五年）：旱灾、蝗灾、疫病、河患与赈济的本年材料连接环境和社会压力，死亡与流离数字须谨慎处理。（材料核查重点：报称信息的来源、抄传与行政分类）}] 明与各地社会、边疆政权；朝廷与官员；跨区域行政、资源或军政网络；不假定同步执行。自然风险与地方报告促成朝廷或地方的应对记录。 赈济、迁徙和损失的实际范围仍须以受灾地区材料分别核验。 材料：《崇祯长编》崇祯五年条及《明史·思宗本纪》；《明史·五行志》；受灾府县地方志与奏疏。限度：桥接条目只标示可回查的年度问题与材料链；实录有中央选择性，崇祯期尤其需要奏疏、地方、朝鲜及满文材料交叉。；本条特别核对报称信息的来源、抄传与行政分类。
  \item[\eventanchor{ML-1632-007}{1632（崇祯五年）：辽饷、剿饷、练饷及地方征解的本年文书显示财政负担，定额、征收与实际流向不得混为一谈。（材料核查重点：报称信息的来源、抄传与行政分类）}] 明与各地社会、边疆政权；朝廷与官员、军队与卫所、地方社会与精英、财政与商业行动者；地方或省域；未具城镇者不补造路线。财政征解、交通工程或货币支付的压力使相关措施进入政务。 其法定安排不等于地方执行，后续须以地域材料追踪负担与成效。 材料：《崇祯长编》崇祯五年条及《明史·思宗本纪》；《明会典》相关条目；地方志、黄册/契约/河工材料。限度：桥接条目只标示可回查的年度问题与材料链；实录有中央选择性，崇祯期尤其需要奏疏、地方、朝鲜及满文材料交叉。；本条特别核对报称信息的来源、抄传与行政分类。
  \item[\eventanchor{ML-1633-001}{1633（崇祯六年）四月：吴桥兵变：山东叛军投奔后金}] 明、后金（清）与辽东诸势力；朝廷与官员、军队与卫所、边地政权与社群、地方社会与精英；边地接触区；不把族名或部名当固定疆界。从边地接触区看「吴桥兵变：山东叛军投奔后金」，朝廷与官员、军队与卫所、边地政权与社群、地方社会与精英的行动连接着当年的军政网络；资料没有充分说明何种先行安排直接促成它，不能把编年相邻当作因果证明。 作为后续影响，「吴桥兵变：山东叛军投奔后金」改变了朝廷与官员、军队与卫所、边地政权与社群、地方社会与精英在边地接触区继续调度、驻守或交涉的起点；能确认的是条件变化，而不是由单一叙述推算出的完整战果。 材料：《崇祯长编》崇祯六年条及《明史·思宗本纪》；《明史·兵志》及相关列传；边镇、地方志或对方材料。限度：「吴桥兵变：山东叛军投奔后金」可由《崇祯长编》崇祯六年条及《明史·思宗本纪》；《明史·兵志》及相关列传；边镇、地方志或对方材料定位到1633（崇祯六年）四月，但这些叙述未完整保存边地接触区的执行与受影响者经验，不能把沉默写成事实。
  \item[\eventanchor{ML-1633-003}{1633（崇祯六年）夏：围攻旅顺：后金占领旅顺}] 明、后金（清）与辽东诸势力；朝廷与官员、军队与卫所、边地政权与社群；边地接触区；不把族名或部名当固定疆界。对于明、后金（清）与辽东诸势力的「围攻旅顺：后金占领旅顺」，边地接触区与朝廷与官员、军队与卫所、边地政权与社群限定了行动的可行范围；史料未逐项记录前置命令和供给，不能以后来结果补造一个整齐的起因。 「围攻旅顺：后金占领旅顺」之后，朝廷与官员、军队与卫所、边地政权与社群在边地接触区的行动空间有所改写；这种变化须由随后文书和地方材料追踪，不能把初报数字直接当作最终后果。 材料：《崇祯长编》崇祯六年条及《明史·思宗本纪》；《明史·兵志》及相关列传；边镇、地方志或对方材料。限度：关于「围攻旅顺：后金占领旅顺」，《崇祯长编》崇祯六年条及《明史·思宗本纪》；《明史·兵志》及相关列传；边镇、地方志或对方材料所能确证的是条目所示的时点和行动；边地接触区内朝廷与官员、军队与卫所、边地政权与社群的规模、意图及社会后果需要异源材料才能判断。
  \item[\eventanchor{ML-1633-004}{1633（崇祯六年）十月22日：辽罗湾海战：明军在金门岛附近击败荷兰海盗船队}] 明与海上、朝贡相关政权；朝廷与官员、军队与卫所；账本可辨空间；地点细节仍须回到所列材料。朝贡名义、边境安全与港口贸易的利益使交涉持续发生。 它为后续册封、互市或海防安排留下文书与跨政权比较线索。 材料：《崇祯长编》崇祯六年条及《明史·思宗本纪》；《明史·外国传》；《朝鲜王朝实录》、琉球《历代宝案》或海外档案。限度：《崇祯长编》崇祯六年条及《明史·思宗本纪》；《明史·外国传》；《朝鲜王朝实录》、琉球《历代宝案》或海外档案使「辽罗湾海战：明军在金门岛附近击败荷兰海盗船队」的时间位置可追查；至于朝廷与官员、军队与卫所在账本可辨空间的命令传递、路线和损失，必须保持待核而不能随编年补写。
  \item[\eventanchor{ML-1633-005}{1633（崇祯六年）十二月27日：叛军占领绵芝}] 明与地方反抗、流动作战集团；朝廷与官员、军队与卫所、地方社会与精英、流动与反抗集团；地方或省域；未具城镇者不补造路线。继承、官僚协调或皇权运作的压力使问题进入朝廷程序。 它影响后续人事与政策议程；动机和派系标签不能由单一叙事裁定。 材料：《崇祯长编》崇祯六年条及《明史·思宗本纪》；《明史·本纪》及相关列传；诏令、奏疏或会典版本。限度：《崇祯长编》崇祯六年条及《明史·思宗本纪》；《明史·本纪》及相关列传；诏令、奏疏或会典版本为「叛军占领绵芝」留下编年入口，却未必保留地方或省域中朝廷与官员、军队与卫所、地方社会与精英、流动与反抗集团的完整视角；本条因此不据之裁定人数、责任或动机。
  \item[\eventanchor{ML-1633-006}{1633（崇祯六年）：天启末至崇祯朝的任免、奏疏和廷议构成本年中枢危机线索；崇祯无实录，必须以多类材料互校。}] 明与各地社会、边疆政权；朝廷与官员；跨区域行政、资源或军政网络；不假定同步执行。继承、官僚协调或皇权运作的压力使问题进入朝廷程序。 它影响后续人事与政策议程；动机和派系标签不能由单一叙事裁定。 材料：《崇祯长编》崇祯六年条及《明史·思宗本纪》；《明史·本纪》及相关列传；诏令、奏疏或会典版本。限度：桥接条目只标示可回查的年度问题与材料链；实录有中央选择性，崇祯期尤其需要奏疏、地方、朝鲜及满文材料交叉。
  \item[\eventanchor{ML-1633-007}{1633（崇祯六年）：辽东防务、后金（清）军事行动和流动作战集团的本年记录标记多战场互动，军数和胜败须保留分歧。}] 明与各地社会、边疆政权；朝廷与官员、军队与卫所、边地政权与社群、流动与反抗集团；边地接触区；不把族名或部名当固定疆界。「辽东防务、后金（清）军事行动和流动作战集团的本年记录标记多战场互动，军数和胜败须保留分歧」把朝廷与官员、军队与卫所、边地政权与社群、流动与反抗集团带到边地接触区的具体关口；现存年条只留下这一行动节点，前期兵力、粮道和地方协力的次序未逐项保存，不能倒推为单一动因。 「辽东防务、后金（清）军事行动和流动作战集团的本年记录标记多战场互动，军数和胜败须保留分歧」改变了朝廷与官员、军队与卫所、边地政权与社群、流动与反抗集团处理边地接触区事务的顺序与代价；后续影响须由连贯记录显示，不能从孤立军报推出可量化的整体结局。 材料：《崇祯长编》崇祯六年条及《明史·思宗本纪》；《明史·兵志》及相关列传；边镇、地方志或对方材料。限度：桥接条目只标示可回查的年度问题与材料链；实录有中央选择性，崇祯期尤其需要奏疏、地方、朝鲜及满文材料交叉。
  \item[\eventanchor{ML-1634-001}{1634（崇祯七年）九月14日：桐城叛乱爆发}] 明与地方反抗、流动作战集团；朝廷与官员、军队与卫所、地方社会与精英、流动与反抗集团；地方或省域；未具城镇者不补造路线。「桐城叛乱爆发」使朝廷与官员、军队与卫所、地方社会与精英、流动与反抗集团在地方或省域的处境变得可见；本条未保存从征发到出动的全链条，故以可核条件而非概括性“压力”标示其背景。 发生「桐城叛乱爆发」后，地方或省域仍是朝廷与官员、军队与卫所、地方社会与精英、流动与反抗集团必须面对的行动场域；可见结果是问题被带入下一段年序，未见的伤亡、掠夺与居民去留不在本条擅断。 材料：《崇祯长编》崇祯七年条及《明史·思宗本纪》；《明史·兵志》及相关列传；边镇、地方志或对方材料。限度：《崇祯长编》崇祯七年条及《明史·思宗本纪》；《明史·兵志》及相关列传；边镇、地方志或对方材料记录「桐城叛乱爆发」的方式限定了可说范围：可追踪1634（崇祯七年）九月14日的行动节点，不能凭此还原地方或省域中朝廷与官员、军队与卫所、地方社会与精英、流动与反抗集团的每一处部署。
  \item[\eventanchor{ML-1634-002}{1634（崇祯七年）：地方结社、宗教、出版或士人文集的本年线索提供战乱中的社会观察，幸存者记忆不能概括全国。}] 明与各地社会、边疆政权；朝廷与官员、军队与卫所、地方社会与精英、宗教与文化行动者；地方或省域；未具城镇者不补造路线。制度、教育或知识传播的需要推动了相关行动或文本形成。 它提供制度和传播史的锚点，不能据此推断社会整体接受。 材料：《崇祯长编》崇祯七年条及《明史·思宗本纪》；《明史·选举志》或《艺文志》；文集、碑刻或书目材料。限度：桥接条目只标示可回查的年度问题与材料链；实录有中央选择性，崇祯期尤其需要奏疏、地方、朝鲜及满文材料交叉。
  \item[\eventanchor{ML-1634-003}{1634（崇祯七年）：辽饷、剿饷、练饷及地方征解的本年文书显示财政负担，定额、征收与实际流向不得混为一谈。（材料核查重点：诏令或奏议的形成与发受链）}] 明与各地社会、边疆政权；朝廷与官员、军队与卫所、地方社会与精英、财政与商业行动者；地方或省域；未具城镇者不补造路线。财政征解、交通工程或货币支付的压力使相关措施进入政务。 其法定安排不等于地方执行，后续须以地域材料追踪负担与成效。 材料：《崇祯长编》崇祯七年条及《明史·思宗本纪》；《明会典》相关条目；地方志、黄册/契约/河工材料。限度：桥接条目只标示可回查的年度问题与材料链；实录有中央选择性，崇祯期尤其需要奏疏、地方、朝鲜及满文材料交叉。；本条特别核对诏令或奏议的形成与发受链。
  \item[\eventanchor{ML-1634-004}{1634（崇祯七年）：辽东防务、后金（清）军事行动和流动作战集团的本年记录标记多战场互动，军数和胜败须保留分歧。}] 明与各地社会、边疆政权；朝廷与官员、军队与卫所、边地政权与社群、流动与反抗集团；边地接触区；不把族名或部名当固定疆界。就「辽东防务、后金（清）军事行动和流动作战集团的本年记录标记多战场互动，军数和胜败须保留分歧」而言，明与各地社会、边疆政权与朝廷与官员、军队与卫所、边地政权与社群、流动与反抗集团在边地接触区的行动并非凭空发生；不过现存材料只标示节点，前序命令、运输和地方响应仍须逐件核验。 「辽东防务、后金（清）军事行动和流动作战集团的本年记录标记多战场互动，军数和胜败须保留分歧」改变了朝廷与官员、军队与卫所、边地政权与社群、流动与反抗集团处理边地接触区事务的顺序与代价；后续影响须由连贯记录显示，不能从孤立军报推出可量化的整体结局。 材料：《崇祯长编》崇祯七年条及《明史·思宗本纪》；《明史·兵志》及相关列传；边镇、地方志或对方材料。限度：桥接条目只标示可回查的年度问题与材料链；实录有中央选择性，崇祯期尤其需要奏疏、地方、朝鲜及满文材料交叉。
  \item[\eventanchor{ML-1634-005}{1634（崇祯七年）：天启末至崇祯朝的任免、奏疏和廷议构成本年中枢危机线索；崇祯无实录，必须以多类材料互校。}] 明与各地社会、边疆政权；朝廷与官员；跨区域行政、资源或军政网络；不假定同步执行。继承、官僚协调或皇权运作的压力使问题进入朝廷程序。 它影响后续人事与政策议程；动机和派系标签不能由单一叙事裁定。 材料：《崇祯长编》崇祯七年条及《明史·思宗本纪》；《明史·本纪》及相关列传；诏令、奏疏或会典版本。限度：桥接条目只标示可回查的年度问题与材料链；实录有中央选择性，崇祯期尤其需要奏疏、地方、朝鲜及满文材料交叉。
  \item[\eventanchor{ML-1634-006}{1634（崇祯七年）：辽饷、剿饷、练饷及地方征解的本年文书显示财政负担，定额、征收与实际流向不得混为一谈。（材料核查重点：报称信息的来源、抄传与行政分类）}] 明与各地社会、边疆政权；朝廷与官员、军队与卫所、地方社会与精英、财政与商业行动者；地方或省域；未具城镇者不补造路线。财政征解、交通工程或货币支付的压力使相关措施进入政务。 其法定安排不等于地方执行，后续须以地域材料追踪负担与成效。 材料：《崇祯长编》崇祯七年条及《明史·思宗本纪》；《明会典》相关条目；地方志、黄册/契约/河工材料。限度：桥接条目只标示可回查的年度问题与材料链；实录有中央选择性，崇祯期尤其需要奏疏、地方、朝鲜及满文材料交叉。；本条特别核对报称信息的来源、抄传与行政分类。
  \item[\eventanchor{ML-1634-007}{1634（崇祯七年）：朝鲜、辽东和海上关系的本年材料提供外部观察位置，明、后金（清）与地方势力的称谓需具体化。}] 明与各地社会、边疆政权；朝廷与官员、边地政权与社群、地方社会与精英；账本可辨空间；地点细节仍须回到所列材料。朝贡名义、边境安全与港口贸易的利益使交涉持续发生。 它为后续册封、互市或海防安排留下文书与跨政权比较线索。 材料：《崇祯长编》崇祯七年条及《明史·思宗本纪》；《明史·外国传》；《朝鲜王朝实录》、琉球《历代宝案》或海外档案。限度：桥接条目只标示可回查的年度问题与材料链；实录有中央选择性，崇祯期尤其需要奏疏、地方、朝鲜及满文材料交叉。
  \item[\eventanchor{ML-1635-001}{1635（崇祯八年）三月：叛军占领凤阳}] 明与地方反抗、流动作战集团；朝廷与官员、军队与卫所、地方社会与精英、流动与反抗集团；地方或省域；未具城镇者不补造路线。继承、官僚协调或皇权运作的压力使问题进入朝廷程序。 它影响后续人事与政策议程；动机和派系标签不能由单一叙事裁定。 材料：《崇祯长编》崇祯八年条及《明史·思宗本纪》；《明史·本纪》及相关列传；诏令、奏疏或会典版本。限度：「叛军占领凤阳」虽见于《崇祯长编》崇祯八年条及《明史·思宗本纪》；《明史·本纪》及相关列传；诏令、奏疏或会典版本，但材料形成者未必观察到地方或省域的全部过程；朝廷与官员、军队与卫所、地方社会与精英、流动与反抗集团的行动范围和受影响群体不能由此单独决定。
  \item[\eventanchor{ML-1635-002}{1635（崇祯八年）八月：明军在甘肃被叛军击败}] 明与地方反抗、流动作战集团；朝廷与官员、军队与卫所、地方社会与精英、流动与反抗集团；地方或省域；未具城镇者不补造路线。「明军在甘肃被叛军击败」出现时，朝廷与官员、军队与卫所、地方社会与精英、流动与反抗集团是在地方或省域的限制下行动；资料所给的是年序定位而非动机自述，前期军政压力只能作为解释线索，不能充作定论。 「明军在甘肃被叛军击败」之后，地方或省域内朝廷与官员、军队与卫所、地方社会与精英、流动与反抗集团仍须处理由该行动留下的控制、通行或防务问题；本条不把战报中的胜负、斩获或伤亡扩大为无从核对的结论。 材料：《崇祯长编》崇祯八年条及《明史·思宗本纪》；《明史·兵志》及相关列传；边镇、地方志或对方材料。限度：《崇祯长编》崇祯八年条及《明史·思宗本纪》；《明史·兵志》及相关列传；边镇、地方志或对方材料使「明军在甘肃被叛军击败」的时间位置可追查；至于朝廷与官员、军队与卫所、地方社会与精英、流动与反抗集团在地方或省域的命令传递、路线和损失，必须保持待核而不能随编年补写。
  \item[\eventanchor{ML-1635-003}{1635（崇祯八年）九月：李自成山西起义军}] 明与地方反抗、流动作战集团；朝廷与官员、军队与卫所、地方社会与精英、流动与反抗集团；地方或省域；未具城镇者不补造路线。继承、官僚协调或皇权运作的压力使问题进入朝廷程序。 它影响后续人事与政策议程；动机和派系标签不能由单一叙事裁定。 材料：《崇祯长编》崇祯八年条及《明史·思宗本纪》；《明史·本纪》及相关列传；诏令、奏疏或会典版本。限度：《崇祯长编》崇祯八年条及《明史·思宗本纪》；《明史·本纪》及相关列传；诏令、奏疏或会典版本为「李自成山西起义军」留下编年入口，却未必保留地方或省域中朝廷与官员、军队与卫所、地方社会与精英、流动与反抗集团的完整视角；本条因此不据之裁定人数、责任或动机。
  \item[\eventanchor{ML-1635-004}{1635（崇祯八年）：明代使用望远镜来瞄准火炮。}] 明与各地社会、边疆政权；朝廷与官员；跨区域行政、资源或军政网络；不假定同步执行。制度、教育或知识传播的需要推动了相关行动或文本形成。 它提供制度和传播史的锚点，不能据此推断社会整体接受。 材料：《崇祯长编》崇祯八年条及《明史·思宗本纪》；《明史·选举志》或《艺文志》；文集、碑刻或书目材料。限度：《崇祯长编》崇祯八年条及《明史·思宗本纪》；《明史·选举志》或《艺文志》；文集、碑刻或书目材料记录「明代使用望远镜来瞄准火炮」的方式限定了可说范围：可追踪1635（崇祯八年）的行动节点，不能凭此还原跨区域行政、资源或军政网络中朝廷与官员的每一处部署。
  \item[\eventanchor{ML-1635-005}{1635（崇祯八年）：辽饷、剿饷、练饷及地方征解的本年文书显示财政负担，定额、征收与实际流向不得混为一谈。}] 明与各地社会、边疆政权；朝廷与官员、军队与卫所、地方社会与精英、财政与商业行动者；地方或省域；未具城镇者不补造路线。财政征解、交通工程或货币支付的压力使相关措施进入政务。 其法定安排不等于地方执行，后续须以地域材料追踪负担与成效。 材料：《崇祯长编》崇祯八年条及《明史·思宗本纪》；《明会典》相关条目；地方志、黄册/契约/河工材料。限度：桥接条目只标示可回查的年度问题与材料链；实录有中央选择性，崇祯期尤其需要奏疏、地方、朝鲜及满文材料交叉。
  \item[\eventanchor{ML-1635-007}{1635（崇祯八年）：地方结社、宗教、出版或士人文集的本年线索提供战乱中的社会观察，幸存者记忆不能概括全国。}] 明与各地社会、边疆政权；朝廷与官员、军队与卫所、地方社会与精英、宗教与文化行动者；地方或省域；未具城镇者不补造路线。制度、教育或知识传播的需要推动了相关行动或文本形成。 它提供制度和传播史的锚点，不能据此推断社会整体接受。 材料：《崇祯长编》崇祯八年条及《明史·思宗本纪》；《明史·选举志》或《艺文志》；文集、碑刻或书目材料。限度：桥接条目只标示可回查的年度问题与材料链；实录有中央选择性，崇祯期尤其需要奏疏、地方、朝鲜及满文材料交叉。
  \item[\eventanchor{ML-1636-001}{1636（崇祯九年）：皇太极宣布清朝成立}] 明、后金（清）与辽东诸势力；朝廷与官员、边地政权与社群；边地接触区；不把族名或部名当固定疆界。继承、官僚协调或皇权运作的压力使问题进入朝廷程序。 它影响后续人事与政策议程；动机和派系标签不能由单一叙事裁定。 材料：《崇祯长编》崇祯九年条及《明史·思宗本纪》；《明史·本纪》及相关列传；诏令、奏疏或会典版本。限度：以《崇祯长编》崇祯九年条及《明史·思宗本纪》；《明史·本纪》及相关列传；诏令、奏疏或会典版本核对「皇太极宣布清朝成立」时，只能确认其在1636（崇祯九年）的叙述位置；朝廷与官员、边地政权与社群在边地接触区的具体经验仍需由可定位的异类材料补证。
  \item[\eventanchor{ML-1636-002}{1636（崇祯九年）：旱灾、蝗灾、疫病、河患与赈济的本年材料连接环境和社会压力，死亡与流离数字须谨慎处理。（材料核查重点：诏令或奏议的形成与发受链）}] 明与各地社会、边疆政权；朝廷与官员；跨区域行政、资源或军政网络；不假定同步执行。自然风险与地方报告促成朝廷或地方的应对记录。 赈济、迁徙和损失的实际范围仍须以受灾地区材料分别核验。 材料：《崇祯长编》崇祯九年条及《明史·思宗本纪》；《明史·五行志》；受灾府县地方志与奏疏。限度：桥接条目只标示可回查的年度问题与材料链；实录有中央选择性，崇祯期尤其需要奏疏、地方、朝鲜及满文材料交叉。；本条特别核对诏令或奏议的形成与发受链。
  \item[\eventanchor{ML-1636-003}{1636（崇祯九年）：辽东防务、后金（清）军事行动和流动作战集团的本年记录标记多战场互动，军数和胜败须保留分歧。（材料核查重点：诏令或奏议的形成与发受链）}] 明与各地社会、边疆政权；朝廷与官员、军队与卫所、边地政权与社群、流动与反抗集团；边地接触区；不把族名或部名当固定疆界。围绕「辽东防务、后金（清）军事行动和流动作战集团的本年记录标记多战场互动，军数和胜败须保留分歧」，明与各地社会、边疆政权在边地接触区面对的是朝廷与官员、军队与卫所、边地政权与社群、流动与反抗集团的现场处置；材料确认行动进入编年，却不足以把此前调兵、补给或地方压力串成确定因果链。 发生「辽东防务、后金（清）军事行动和流动作战集团的本年记录标记多战场互动，军数和胜败须保留分歧」后，边地接触区仍是朝廷与官员、军队与卫所、边地政权与社群、流动与反抗集团必须面对的行动场域；可见结果是问题被带入下一段年序，未见的伤亡、掠夺与居民去留不在本条擅断。 材料：《崇祯长编》崇祯九年条及《明史·思宗本纪》；《明史·兵志》及相关列传；边镇、地方志或对方材料。限度：桥接条目只标示可回查的年度问题与材料链；实录有中央选择性，崇祯期尤其需要奏疏、地方、朝鲜及满文材料交叉。；本条特别核对诏令或奏议的形成与发受链。
  \item[\eventanchor{ML-1636-004}{1636（崇祯九年）：朝鲜、辽东和海上关系的本年材料提供外部观察位置，明、后金（清）与地方势力的称谓需具体化。（材料核查重点：诏令或奏议的形成与发受链）}] 明与各地社会、边疆政权；朝廷与官员、边地政权与社群、地方社会与精英；账本可辨空间；地点细节仍须回到所列材料。朝贡名义、边境安全与港口贸易的利益使交涉持续发生。 它为后续册封、互市或海防安排留下文书与跨政权比较线索。 材料：《崇祯长编》崇祯九年条及《明史·思宗本纪》；《明史·外国传》；《朝鲜王朝实录》、琉球《历代宝案》或海外档案。限度：桥接条目只标示可回查的年度问题与材料链；实录有中央选择性，崇祯期尤其需要奏疏、地方、朝鲜及满文材料交叉。；本条特别核对诏令或奏议的形成与发受链。
  \item[\eventanchor{ML-1636-005}{1636（崇祯九年）：辽东防务、后金（清）军事行动和流动作战集团的本年记录标记多战场互动，军数和胜败须保留分歧。（材料核查重点：报称信息的来源、抄传与行政分类）}] 明与各地社会、边疆政权；朝廷与官员、军队与卫所、边地政权与社群、流动与反抗集团；边地接触区；不把族名或部名当固定疆界。「辽东防务、后金（清）军事行动和流动作战集团的本年记录标记多战场互动，军数和胜败须保留分歧」记录了朝廷与官员、军队与卫所、边地政权与社群、流动与反抗集团在边地接触区的一次处置；其兵员、道路或地方支持并非都见于同一材料，因而不以总括理由覆盖未见环节。 「辽东防务、后金（清）军事行动和流动作战集团的本年记录标记多战场互动，军数和胜败须保留分歧」之后，边地接触区内朝廷与官员、军队与卫所、边地政权与社群、流动与反抗集团仍须处理由该行动留下的控制、通行或防务问题；本条不把战报中的胜负、斩获或伤亡扩大为无从核对的结论。 材料：《崇祯长编》崇祯九年条及《明史·思宗本纪》；《明史·兵志》及相关列传；边镇、地方志或对方材料。限度：桥接条目只标示可回查的年度问题与材料链；实录有中央选择性，崇祯期尤其需要奏疏、地方、朝鲜及满文材料交叉。；本条特别核对报称信息的来源、抄传与行政分类。
  \item[\eventanchor{ML-1636-006}{1636（崇祯九年）：朝鲜、辽东和海上关系的本年材料提供外部观察位置，明、后金（清）与地方势力的称谓需具体化。（材料核查重点：报称信息的来源、抄传与行政分类）}] 明与各地社会、边疆政权；朝廷与官员、边地政权与社群、地方社会与精英；账本可辨空间；地点细节仍须回到所列材料。朝贡名义、边境安全与港口贸易的利益使交涉持续发生。 它为后续册封、互市或海防安排留下文书与跨政权比较线索。 材料：《崇祯长编》崇祯九年条及《明史·思宗本纪》；《明史·外国传》；《朝鲜王朝实录》、琉球《历代宝案》或海外档案。限度：桥接条目只标示可回查的年度问题与材料链；实录有中央选择性，崇祯期尤其需要奏疏、地方、朝鲜及满文材料交叉。；本条特别核对报称信息的来源、抄传与行政分类。
  \item[\eventanchor{ML-1636-007}{1636（崇祯九年）：旱灾、蝗灾、疫病、河患与赈济的本年材料连接环境和社会压力，死亡与流离数字须谨慎处理。（材料核查重点：报称信息的来源、抄传与行政分类）}] 明与各地社会、边疆政权；朝廷与官员；跨区域行政、资源或军政网络；不假定同步执行。自然风险与地方报告促成朝廷或地方的应对记录。 赈济、迁徙和损失的实际范围仍须以受灾地区材料分别核验。 材料：《崇祯长编》崇祯九年条及《明史·思宗本纪》；《明史·五行志》；受灾府县地方志与奏疏。限度：桥接条目只标示可回查的年度问题与材料链；实录有中央选择性，崇祯期尤其需要奏疏、地方、朝鲜及满文材料交叉。；本条特别核对报称信息的来源、抄传与行政分类。
  \item[\eventanchor{ML-1637-001}{1637（崇祯十年）：天启末至崇祯朝的任免、奏疏和廷议构成本年中枢危机线索；崇祯无实录，必须以多类材料互校。（材料核查重点：诏令或奏议的形成与发受链）}] 明与各地社会、边疆政权；朝廷与官员；跨区域行政、资源或军政网络；不假定同步执行。继承、官僚协调或皇权运作的压力使问题进入朝廷程序。 它影响后续人事与政策议程；动机和派系标签不能由单一叙事裁定。 材料：《崇祯长编》崇祯十年条及《明史·思宗本纪》；《明史·本纪》及相关列传；诏令、奏疏或会典版本。限度：桥接条目只标示可回查的年度问题与材料链；实录有中央选择性，崇祯期尤其需要奏疏、地方、朝鲜及满文材料交叉。；本条特别核对诏令或奏议的形成与发受链。
  \item[\eventanchor{ML-1637-002}{1637（崇祯十年）：天启末至崇祯朝的任免、奏疏和廷议构成本年中枢危机线索；崇祯无实录，必须以多类材料互校。（材料核查重点：报称信息的来源、抄传与行政分类）}] 明与各地社会、边疆政权；朝廷与官员；跨区域行政、资源或军政网络；不假定同步执行。继承、官僚协调或皇权运作的压力使问题进入朝廷程序。 它影响后续人事与政策议程；动机和派系标签不能由单一叙事裁定。 材料：《崇祯长编》崇祯十年条及《明史·思宗本纪》；《明史·本纪》及相关列传；诏令、奏疏或会典版本。限度：桥接条目只标示可回查的年度问题与材料链；实录有中央选择性，崇祯期尤其需要奏疏、地方、朝鲜及满文材料交叉。；本条特别核对报称信息的来源、抄传与行政分类。
  \item[\eventanchor{ML-1637-003}{1637（崇祯十年）：朝鲜、辽东和海上关系的本年材料提供外部观察位置，明、后金（清）与地方势力的称谓需具体化。（材料核查重点：诏令或奏议的形成与发受链）}] 明与各地社会、边疆政权；朝廷与官员、边地政权与社群、地方社会与精英；账本可辨空间；地点细节仍须回到所列材料。朝贡名义、边境安全与港口贸易的利益使交涉持续发生。 它为后续册封、互市或海防安排留下文书与跨政权比较线索。 材料：《崇祯长编》崇祯十年条及《明史·思宗本纪》；《明史·外国传》；《朝鲜王朝实录》、琉球《历代宝案》或海外档案。限度：桥接条目只标示可回查的年度问题与材料链；实录有中央选择性，崇祯期尤其需要奏疏、地方、朝鲜及满文材料交叉。；本条特别核对诏令或奏议的形成与发受链。
  \item[\eventanchor{ML-1637-004}{1637（崇祯十年）：地方结社、宗教、出版或士人文集的本年线索提供战乱中的社会观察，幸存者记忆不能概括全国。（材料核查重点：诏令或奏议的形成与发受链）}] 明与各地社会、边疆政权；朝廷与官员、军队与卫所、地方社会与精英、宗教与文化行动者；地方或省域；未具城镇者不补造路线。制度、教育或知识传播的需要推动了相关行动或文本形成。 它提供制度和传播史的锚点，不能据此推断社会整体接受。 材料：《崇祯长编》崇祯十年条及《明史·思宗本纪》；《明史·选举志》或《艺文志》；文集、碑刻或书目材料。限度：桥接条目只标示可回查的年度问题与材料链；实录有中央选择性，崇祯期尤其需要奏疏、地方、朝鲜及满文材料交叉。；本条特别核对诏令或奏议的形成与发受链。
  \item[\eventanchor{ML-1637-005}{1637（崇祯十年）：朝鲜、辽东和海上关系的本年材料提供外部观察位置，明、后金（清）与地方势力的称谓需具体化。（材料核查重点：报称信息的来源、抄传与行政分类）}] 明与各地社会、边疆政权；朝廷与官员、边地政权与社群、地方社会与精英；账本可辨空间；地点细节仍须回到所列材料。朝贡名义、边境安全与港口贸易的利益使交涉持续发生。 它为后续册封、互市或海防安排留下文书与跨政权比较线索。 材料：《崇祯长编》崇祯十年条及《明史·思宗本纪》；《明史·外国传》；《朝鲜王朝实录》、琉球《历代宝案》或海外档案。限度：桥接条目只标示可回查的年度问题与材料链；实录有中央选择性，崇祯期尤其需要奏疏、地方、朝鲜及满文材料交叉。；本条特别核对报称信息的来源、抄传与行政分类。
  \item[\eventanchor{ML-1637-006}{1637（崇祯十年）：地方结社、宗教、出版或士人文集的本年线索提供战乱中的社会观察，幸存者记忆不能概括全国。（材料核查重点：报称信息的来源、抄传与行政分类）}] 明与各地社会、边疆政权；朝廷与官员、军队与卫所、地方社会与精英、宗教与文化行动者；地方或省域；未具城镇者不补造路线。制度、教育或知识传播的需要推动了相关行动或文本形成。 它提供制度和传播史的锚点，不能据此推断社会整体接受。 材料：《崇祯长编》崇祯十年条及《明史·思宗本纪》；《明史·选举志》或《艺文志》；文集、碑刻或书目材料。限度：桥接条目只标示可回查的年度问题与材料链；实录有中央选择性，崇祯期尤其需要奏疏、地方、朝鲜及满文材料交叉。；本条特别核对报称信息的来源、抄传与行政分类。
  \item[\eventanchor{ML-1637-007}{1637（崇祯十年）：天启末至崇祯朝的任免、奏疏和廷议构成本年中枢危机线索；崇祯无实录，必须以多类材料互校。（材料核查重点：同年地方材料所见的执行范围）}] 明与各地社会、边疆政权；朝廷与官员、军队与卫所、地方社会与精英；地方或省域；未具城镇者不补造路线。继承、官僚协调或皇权运作的压力使问题进入朝廷程序。 它影响后续人事与政策议程；动机和派系标签不能由单一叙事裁定。 材料：《崇祯长编》崇祯十年条及《明史·思宗本纪》；《明史·本纪》及相关列传；诏令、奏疏或会典版本。限度：桥接条目只标示可回查的年度问题与材料链；实录有中央选择性，崇祯期尤其需要奏疏、地方、朝鲜及满文材料交叉。；本条特别核对同年地方材料所见的执行范围。
  \item[\eventanchor{ML-1638-001}{1638（崇祯十一年）：清朝征服山东}] 明、后金（清）与辽东诸势力；朝廷与官员、军队与卫所、边地政权与社群、地方社会与精英；边地接触区；不把族名或部名当固定疆界。「清朝征服山东」使朝廷与官员、军队与卫所、边地政权与社群、地方社会与精英在边地接触区的处境变得可见；本条未保存从征发到出动的全链条，故以可核条件而非概括性“压力”标示其背景。 「清朝征服山东」使明、后金（清）与辽东诸势力在边地接触区的下一步选择重新围绕朝廷与官员、军队与卫所、边地政权与社群、地方社会与精英展开；其影响应落到随后可见的军令、安置或交涉，而非以未经互校的人数概括。 材料：《崇祯长编》崇祯十一年条及《明史·思宗本纪》；《明史·兵志》及相关列传；边镇、地方志或对方材料。限度：就「清朝征服山东」而论，《崇祯长编》崇祯十一年条及《明史·思宗本纪》；《明史·兵志》及相关列传；边镇、地方志或对方材料提供的是1638（崇祯十一年）的记录线索；朝廷与官员、军队与卫所、边地政权与社群、地方社会与精英在边地接触区的实际行动细节不能由这一材料链独自补足。
  \item[\eventanchor{ML-1638-002}{1638（崇祯十一年）：明军在晋豫边境大败}] 明与各地社会、边疆政权；朝廷与官员、军队与卫所；跨区域行政、资源或军政网络；不假定同步执行。从跨区域行政、资源或军政网络看「明军在晋豫边境大败」，朝廷与官员、军队与卫所的行动连接着当年的军政网络；资料没有充分说明何种先行安排直接促成它，不能把编年相邻当作因果证明。 「明军在晋豫边境大败」并未结束朝廷与官员、军队与卫所在跨区域行政、资源或军政网络的压力，而是把问题转到后续防务、运输或行政处置；只有后续记录能说明它实际造成何种改变。 材料：《崇祯长编》崇祯十一年条及《明史·思宗本纪》；《明史·兵志》及相关列传；边镇、地方志或对方材料。限度：《崇祯长编》崇祯十一年条及《明史·思宗本纪》；《明史·兵志》及相关列传；边镇、地方志或对方材料为「明军在晋豫边境大败」留下编年入口，却未必保留跨区域行政、资源或军政网络中朝廷与官员、军队与卫所的完整视角；本条因此不据之裁定人数、责任或动机。
  \item[\eventanchor{ML-1638-004}{1638（崇祯十一年）：旱灾、蝗灾、疫病、河患与赈济的本年材料连接环境和社会压力，死亡与流离数字须谨慎处理。（材料核查重点：诏令或奏议的形成与发受链）}] 明与各地社会、边疆政权；朝廷与官员；跨区域行政、资源或军政网络；不假定同步执行。自然风险与地方报告促成朝廷或地方的应对记录。 赈济、迁徙和损失的实际范围仍须以受灾地区材料分别核验。 材料：《崇祯长编》崇祯十一年条及《明史·思宗本纪》；《明史·五行志》；受灾府县地方志与奏疏。限度：桥接条目只标示可回查的年度问题与材料链；实录有中央选择性，崇祯期尤其需要奏疏、地方、朝鲜及满文材料交叉。；本条特别核对诏令或奏议的形成与发受链。
  \item[\eventanchor{ML-1638-005}{1638（崇祯十一年）：地方结社、宗教、出版或士人文集的本年线索提供战乱中的社会观察，幸存者记忆不能概括全国。（材料核查重点：报称信息的来源、抄传与行政分类）}] 明与各地社会、边疆政权；朝廷与官员、军队与卫所、地方社会与精英、宗教与文化行动者；地方或省域；未具城镇者不补造路线。制度、教育或知识传播的需要推动了相关行动或文本形成。 它提供制度和传播史的锚点，不能据此推断社会整体接受。 材料：《崇祯长编》崇祯十一年条及《明史·思宗本纪》；《明史·选举志》或《艺文志》；文集、碑刻或书目材料。限度：桥接条目只标示可回查的年度问题与材料链；实录有中央选择性，崇祯期尤其需要奏疏、地方、朝鲜及满文材料交叉。；本条特别核对报称信息的来源、抄传与行政分类。
  \item[\eventanchor{ML-1638-006}{1638（崇祯十一年）：旱灾、蝗灾、疫病、河患与赈济的本年材料连接环境和社会压力，死亡与流离数字须谨慎处理。（材料核查重点：报称信息的来源、抄传与行政分类）}] 明与各地社会、边疆政权；朝廷与官员；跨区域行政、资源或军政网络；不假定同步执行。自然风险与地方报告促成朝廷或地方的应对记录。 赈济、迁徙和损失的实际范围仍须以受灾地区材料分别核验。 材料：《崇祯长编》崇祯十一年条及《明史·思宗本纪》；《明史·五行志》；受灾府县地方志与奏疏。限度：桥接条目只标示可回查的年度问题与材料链；实录有中央选择性，崇祯期尤其需要奏疏、地方、朝鲜及满文材料交叉。；本条特别核对报称信息的来源、抄传与行政分类。
  \item[\eventanchor{ML-1638-007}{1638（崇祯十一年）：辽饷、剿饷、练饷及地方征解的本年文书显示财政负担，定额、征收与实际流向不得混为一谈。}] 明与各地社会、边疆政权；朝廷与官员、军队与卫所、地方社会与精英、财政与商业行动者；地方或省域；未具城镇者不补造路线。财政征解、交通工程或货币支付的压力使相关措施进入政务。 其法定安排不等于地方执行，后续须以地域材料追踪负担与成效。 材料：《崇祯长编》崇祯十一年条及《明史·思宗本纪》；《明会典》相关条目；地方志、黄册/契约/河工材料。限度：桥接条目只标示可回查的年度问题与材料链；实录有中央选择性，崇祯期尤其需要奏疏、地方、朝鲜及满文材料交叉。
  \item[\eventanchor{ML-1639-001}{1639（崇祯十二年）：西班牙人和菲律宾人在吕宋岛屠杀了两万名中国人}] 明与海上、朝贡相关政权；朝廷与官员；账本可辨空间；地点细节仍须回到所列材料。朝贡名义、边境安全与港口贸易的利益使交涉持续发生。 它为后续册封、互市或海防安排留下文书与跨政权比较线索。 材料：《崇祯长编》崇祯十二年条及《明史·思宗本纪》；《明史·外国传》；《朝鲜王朝实录》、琉球《历代宝案》或海外档案。限度：《崇祯长编》崇祯十二年条及《明史·思宗本纪》；《明史·外国传》；《朝鲜王朝实录》、琉球《历代宝案》或海外档案只能把「西班牙人和菲律宾人在吕宋岛屠杀了两万名中国人」置入1639（崇祯十二年）的年序；对于账本可辨空间中朝廷与官员的路线、人数和动机，仍须另据地方或对方材料核验。
  \item[\eventanchor{ML-1639-002}{1639（崇祯十二年）：禁止来自澳门的葡萄牙商人进入长崎}] 明与海上、朝贡相关政权；朝廷与官员、地方社会与精英、财政与商业行动者；账本可辨空间；地点细节仍须回到所列材料。朝贡名义、边境安全与港口贸易的利益使交涉持续发生。 它为后续册封、互市或海防安排留下文书与跨政权比较线索。 材料：《崇祯长编》崇祯十二年条及《明史·思宗本纪》；《明史·外国传》；《朝鲜王朝实录》、琉球《历代宝案》或海外档案。限度：以《崇祯长编》崇祯十二年条及《明史·思宗本纪》；《明史·外国传》；《朝鲜王朝实录》、琉球《历代宝案》或海外档案核对「禁止来自澳门的葡萄牙商人进入长崎」时，只能确认其在1639（崇祯十二年）的叙述位置；朝廷与官员、地方社会与精英、财政与商业行动者在账本可辨空间的具体经验仍需由可定位的异类材料补证。
  \item[\eventanchor{ML-1639-003}{1639（崇祯十二年）：浙江遭遇旱灾}] 明与各地社会、边疆政权；朝廷与官员；跨区域行政、资源或军政网络；不假定同步执行。自然风险与地方报告促成朝廷或地方的应对记录。 赈济、迁徙和损失的实际范围仍须以受灾地区材料分别核验。 材料：《崇祯长编》崇祯十二年条及《明史·思宗本纪》；《明史·五行志》；受灾府县地方志与奏疏。限度：对「浙江遭遇旱灾」的《崇祯长编》崇祯十二年条及《明史·思宗本纪》；《明史·五行志》；受灾府县地方志与奏疏记录可支持年序核验；它不足以替代跨区域行政、资源或军政网络的地方材料，故不把朝廷与官员的战报、奏报或叙事当成无争议结果。
  \item[\eventanchor{ML-1639-004}{1639（崇祯十二年）：天启末至崇祯朝的任免、奏疏和廷议构成本年中枢危机线索；崇祯无实录，必须以多类材料互校。（材料核查重点：诏令或奏议的形成与发受链）}] 明与各地社会、边疆政权；朝廷与官员；跨区域行政、资源或军政网络；不假定同步执行。继承、官僚协调或皇权运作的压力使问题进入朝廷程序。 它影响后续人事与政策议程；动机和派系标签不能由单一叙事裁定。 材料：《崇祯长编》崇祯十二年条及《明史·思宗本纪》；《明史·本纪》及相关列传；诏令、奏疏或会典版本。限度：桥接条目只标示可回查的年度问题与材料链；实录有中央选择性，崇祯期尤其需要奏疏、地方、朝鲜及满文材料交叉。；本条特别核对诏令或奏议的形成与发受链。
  \item[\eventanchor{ML-1639-005}{1639（崇祯十二年）：旱灾、蝗灾、疫病、河患与赈济的本年材料连接环境和社会压力，死亡与流离数字须谨慎处理。}] 明与各地社会、边疆政权；朝廷与官员；跨区域行政、资源或军政网络；不假定同步执行。自然风险与地方报告促成朝廷或地方的应对记录。 赈济、迁徙和损失的实际范围仍须以受灾地区材料分别核验。 材料：《崇祯长编》崇祯十二年条及《明史·思宗本纪》；《明史·五行志》；受灾府县地方志与奏疏。限度：桥接条目只标示可回查的年度问题与材料链；实录有中央选择性，崇祯期尤其需要奏疏、地方、朝鲜及满文材料交叉。
  \item[\eventanchor{ML-1639-006}{1639（崇祯十二年）：天启末至崇祯朝的任免、奏疏和廷议构成本年中枢危机线索；崇祯无实录，必须以多类材料互校。（材料核查重点：报称信息的来源、抄传与行政分类）}] 明与各地社会、边疆政权；朝廷与官员；跨区域行政、资源或军政网络；不假定同步执行。继承、官僚协调或皇权运作的压力使问题进入朝廷程序。 它影响后续人事与政策议程；动机和派系标签不能由单一叙事裁定。 材料：《崇祯长编》崇祯十二年条及《明史·思宗本纪》；《明史·本纪》及相关列传；诏令、奏疏或会典版本。限度：桥接条目只标示可回查的年度问题与材料链；实录有中央选择性，崇祯期尤其需要奏疏、地方、朝鲜及满文材料交叉。；本条特别核对报称信息的来源、抄传与行政分类。
  \item[\eventanchor{ML-1639-007}{1639（崇祯十二年）：辽东防务、后金（清）军事行动和流动作战集团的本年记录标记多战场互动，军数和胜败须保留分歧。}] 明与各地社会、边疆政权；朝廷与官员、军队与卫所、边地政权与社群、流动与反抗集团；边地接触区；不把族名或部名当固定疆界。从边地接触区看「辽东防务、后金（清）军事行动和流动作战集团的本年记录标记多战场互动，军数和胜败须保留分歧」，朝廷与官员、军队与卫所、边地政权与社群、流动与反抗集团的行动连接着当年的军政网络；资料没有充分说明何种先行安排直接促成它，不能把编年相邻当作因果证明。 作为后续影响，「辽东防务、后金（清）军事行动和流动作战集团的本年记录标记多战场互动，军数和胜败须保留分歧」改变了朝廷与官员、军队与卫所、边地政权与社群、流动与反抗集团在边地接触区继续调度、驻守或交涉的起点；能确认的是条件变化，而不是由单一叙述推算出的完整战果。 材料：《崇祯长编》崇祯十二年条及《明史·思宗本纪》；《明史·兵志》及相关列传；边镇、地方志或对方材料。限度：桥接条目只标示可回查的年度问题与材料链；实录有中央选择性，崇祯期尤其需要奏疏、地方、朝鲜及满文材料交叉。
  \item[\eventanchor{ML-1640-001}{1640（崇祯十三年）夏：叛军进入四川}] 明与地方反抗、流动作战集团；朝廷与官员、军队与卫所、地方社会与精英、流动与反抗集团；地方或省域；未具城镇者不补造路线。继承、官僚协调或皇权运作的压力使问题进入朝廷程序。 它影响后续人事与政策议程；动机和派系标签不能由单一叙事裁定。 材料：《崇祯长编》崇祯十三年条及《明史·思宗本纪》；《明史·本纪》及相关列传；诏令、奏疏或会典版本。限度：「叛军进入四川」可由《崇祯长编》崇祯十三年条及《明史·思宗本纪》；《明史·本纪》及相关列传；诏令、奏疏或会典版本定位到1640（崇祯十三年）夏，但这些叙述未完整保存地方或省域的执行与受影响者经验，不能把沉默写成事实。
  \item[\eventanchor{ML-1640-002}{1640（崇祯十三年）：辽饷、剿饷、练饷及地方征解的本年文书显示财政负担，定额、征收与实际流向不得混为一谈。（材料核查重点：诏令或奏议的形成与发受链）}] 明与各地社会、边疆政权；朝廷与官员、军队与卫所、地方社会与精英、财政与商业行动者；地方或省域；未具城镇者不补造路线。财政征解、交通工程或货币支付的压力使相关措施进入政务。 其法定安排不等于地方执行，后续须以地域材料追踪负担与成效。 材料：《崇祯长编》崇祯十三年条及《明史·思宗本纪》；《明会典》相关条目；地方志、黄册/契约/河工材料。限度：桥接条目只标示可回查的年度问题与材料链；实录有中央选择性，崇祯期尤其需要奏疏、地方、朝鲜及满文材料交叉。；本条特别核对诏令或奏议的形成与发受链。
  \item[\eventanchor{ML-1640-003}{1640（崇祯十三年）：旱灾、蝗灾、疫病、河患与赈济的本年材料连接环境和社会压力，死亡与流离数字须谨慎处理。}] 明与各地社会、边疆政权；朝廷与官员；跨区域行政、资源或军政网络；不假定同步执行。自然风险与地方报告促成朝廷或地方的应对记录。 赈济、迁徙和损失的实际范围仍须以受灾地区材料分别核验。 材料：《崇祯长编》崇祯十三年条及《明史·思宗本纪》；《明史·五行志》；受灾府县地方志与奏疏。限度：桥接条目只标示可回查的年度问题与材料链；实录有中央选择性，崇祯期尤其需要奏疏、地方、朝鲜及满文材料交叉。
  \item[\eventanchor{ML-1640-004}{1640（崇祯十三年）：辽饷、剿饷、练饷及地方征解的本年文书显示财政负担，定额、征收与实际流向不得混为一谈。（材料核查重点：报称信息的来源、抄传与行政分类）}] 明与各地社会、边疆政权；朝廷与官员、军队与卫所、地方社会与精英、财政与商业行动者；地方或省域；未具城镇者不补造路线。财政征解、交通工程或货币支付的压力使相关措施进入政务。 其法定安排不等于地方执行，后续须以地域材料追踪负担与成效。 材料：《崇祯长编》崇祯十三年条及《明史·思宗本纪》；《明会典》相关条目；地方志、黄册/契约/河工材料。限度：桥接条目只标示可回查的年度问题与材料链；实录有中央选择性，崇祯期尤其需要奏疏、地方、朝鲜及满文材料交叉。；本条特别核对报称信息的来源、抄传与行政分类。
  \item[\eventanchor{ML-1640-005}{1640（崇祯十三年）：天启末至崇祯朝的任免、奏疏和廷议构成本年中枢危机线索；崇祯无实录，必须以多类材料互校。}] 明与各地社会、边疆政权；朝廷与官员；跨区域行政、资源或军政网络；不假定同步执行。继承、官僚协调或皇权运作的压力使问题进入朝廷程序。 它影响后续人事与政策议程；动机和派系标签不能由单一叙事裁定。 材料：《崇祯长编》崇祯十三年条及《明史·思宗本纪》；《明史·本纪》及相关列传；诏令、奏疏或会典版本。限度：桥接条目只标示可回查的年度问题与材料链；实录有中央选择性，崇祯期尤其需要奏疏、地方、朝鲜及满文材料交叉。
  \item[\eventanchor{ML-1640-006}{1640（崇祯十三年）：辽饷、剿饷、练饷及地方征解的本年文书显示财政负担，定额、征收与实际流向不得混为一谈。（材料核查重点：同年地方材料所见的执行范围）}] 明与各地社会、边疆政权；朝廷与官员、军队与卫所、地方社会与精英、财政与商业行动者；地方或省域；未具城镇者不补造路线。财政征解、交通工程或货币支付的压力使相关措施进入政务。 其法定安排不等于地方执行，后续须以地域材料追踪负担与成效。 材料：《崇祯长编》崇祯十三年条及《明史·思宗本纪》；《明会典》相关条目；地方志、黄册/契约/河工材料。限度：桥接条目只标示可回查的年度问题与材料链；实录有中央选择性，崇祯期尤其需要奏疏、地方、朝鲜及满文材料交叉。；本条特别核对同年地方材料所见的执行范围。
  \item[\eventanchor{ML-1640-007}{1640（崇祯十三年）：朝鲜、辽东和海上关系的本年材料提供外部观察位置，明、后金（清）与地方势力的称谓需具体化。}] 明与各地社会、边疆政权；朝廷与官员、边地政权与社群、地方社会与精英；账本可辨空间；地点细节仍须回到所列材料。朝贡名义、边境安全与港口贸易的利益使交涉持续发生。 它为后续册封、互市或海防安排留下文书与跨政权比较线索。 材料：《崇祯长编》崇祯十三年条及《明史·思宗本纪》；《明史·外国传》；《朝鲜王朝实录》、琉球《历代宝案》或海外档案。限度：桥接条目只标示可回查的年度问题与材料链；实录有中央选择性，崇祯期尤其需要奏疏、地方、朝鲜及满文材料交叉。
  \item[\eventanchor{ML-1641-002}{1641（崇祯十四年）三月：李自成攻克洛阳}] 明与地方反抗、流动作战集团；朝廷与官员、军队与卫所、地方社会与精英、流动与反抗集团；地方或省域；未具城镇者不补造路线。继承、官僚协调或皇权运作的压力使问题进入朝廷程序。 它影响后续人事与政策议程；动机和派系标签不能由单一叙事裁定。 材料：《崇祯长编》崇祯十四年条及《明史·思宗本纪》；《明史·本纪》及相关列传；诏令、奏疏或会典版本。限度：关于「李自成攻克洛阳」，《崇祯长编》崇祯十四年条及《明史·思宗本纪》；《明史·本纪》及相关列传；诏令、奏疏或会典版本所能确证的是条目所示的时点和行动；地方或省域内朝廷与官员、军队与卫所、地方社会与精英、流动与反抗集团的规模、意图及社会后果需要异源材料才能判断。
  \item[\eventanchor{ML-1641-003}{1641（崇祯十四年）：起义军首领张献忠攻克襄阳}] 明与地方反抗、流动作战集团；朝廷与官员、军队与卫所、地方社会与精英、流动与反抗集团；地方或省域；未具城镇者不补造路线。继承、官僚协调或皇权运作的压力使问题进入朝廷程序。 它影响后续人事与政策议程；动机和派系标签不能由单一叙事裁定。 材料：《崇祯长编》崇祯十四年条及《明史·思宗本纪》；《明史·本纪》及相关列传；诏令、奏疏或会典版本。限度：《崇祯长编》崇祯十四年条及《明史·思宗本纪》；《明史·本纪》及相关列传；诏令、奏疏或会典版本记录「起义军首领张献忠攻克襄阳」的方式限定了可说范围：可追踪1641（崇祯十四年）的行动节点，不能凭此还原地方或省域中朝廷与官员、军队与卫所、地方社会与精英、流动与反抗集团的每一处部署。
  \item[\eventanchor{ML-1641-004}{1641（崇祯十四年）七月15日：张献忠拿下武昌}] 明与地方反抗、流动作战集团；朝廷与官员、军队与卫所、地方社会与精英、流动与反抗集团；地方或省域；未具城镇者不补造路线。继承、官僚协调或皇权运作的压力使问题进入朝廷程序。 它影响后续人事与政策议程；动机和派系标签不能由单一叙事裁定。 材料：《崇祯长编》崇祯十四年条及《明史·思宗本纪》；《明史·本纪》及相关列传；诏令、奏疏或会典版本。限度：「张献忠拿下武昌」可由《崇祯长编》崇祯十四年条及《明史·思宗本纪》；《明史·本纪》及相关列传；诏令、奏疏或会典版本定位到1641（崇祯十四年）七月15日，但这些叙述未完整保存地方或省域的执行与受影响者经验，不能把沉默写成事实。
  \item[\eventanchor{ML-1641-005}{1641（崇祯十四年）：蝗虫袭击浙江}] 明与各地社会、边疆政权；朝廷与官员；跨区域行政、资源或军政网络；不假定同步执行。自然风险与地方报告促成朝廷或地方的应对记录。 赈济、迁徙和损失的实际范围仍须以受灾地区材料分别核验。 材料：《崇祯长编》崇祯十四年条及《明史·思宗本纪》；《明史·五行志》；受灾府县地方志与奏疏。限度：以《崇祯长编》崇祯十四年条及《明史·思宗本纪》；《明史·五行志》；受灾府县地方志与奏疏核对「蝗虫袭击浙江」时，只能确认其在1641（崇祯十四年）的叙述位置；朝廷与官员在跨区域行政、资源或军政网络的具体经验仍需由可定位的异类材料补证。
  \item[\eventanchor{ML-1641-006}{1641（崇祯十四年）：辽东防务、后金（清）军事行动和流动作战集团的本年记录标记多战场互动，军数和胜败须保留分歧。}] 明与各地社会、边疆政权；朝廷与官员、军队与卫所、边地政权与社群、流动与反抗集团；边地接触区；不把族名或部名当固定疆界。就「辽东防务、后金（清）军事行动和流动作战集团的本年记录标记多战场互动，军数和胜败须保留分歧」而言，明与各地社会、边疆政权与朝廷与官员、军队与卫所、边地政权与社群、流动与反抗集团在边地接触区的行动并非凭空发生；不过现存材料只标示节点，前序命令、运输和地方响应仍须逐件核验。 此后，朝廷与官员、军队与卫所、边地政权与社群、流动与反抗集团在边地接触区面对的并非原来的行动条件；「辽东防务、后金（清）军事行动和流动作战集团的本年记录标记多战场互动，军数和胜败须保留分歧」留下需要继续处理的秩序与供给问题，损失和胜负仍应按异源材料分别核对。 材料：《崇祯长编》崇祯十四年条及《明史·思宗本纪》；《明史·兵志》及相关列传；边镇、地方志或对方材料。限度：桥接条目只标示可回查的年度问题与材料链；实录有中央选择性，崇祯期尤其需要奏疏、地方、朝鲜及满文材料交叉。
  \item[\eventanchor{ML-1641-007}{1641（崇祯十四年）：地方结社、宗教、出版或士人文集的本年线索提供战乱中的社会观察，幸存者记忆不能概括全国。}] 明与各地社会、边疆政权；朝廷与官员、军队与卫所、地方社会与精英、宗教与文化行动者；地方或省域；未具城镇者不补造路线。制度、教育或知识传播的需要推动了相关行动或文本形成。 它提供制度和传播史的锚点，不能据此推断社会整体接受。 材料：《崇祯长编》崇祯十四年条及《明史·思宗本纪》；《明史·选举志》或《艺文志》；文集、碑刻或书目材料。限度：桥接条目只标示可回查的年度问题与材料链；实录有中央选择性，崇祯期尤其需要奏疏、地方、朝鲜及满文材料交叉。
  \item[\eventanchor{ML-1642-001}{1642（崇祯十五年）四月8日：宋金之战：清军占领锦州}] 明、后金（清）与辽东诸势力；朝廷与官员、军队与卫所、边地政权与社群；边地接触区；不把族名或部名当固定疆界。「宋金之战：清军占领锦州」把朝廷与官员、军队与卫所、边地政权与社群带到边地接触区的具体关口；现存年条只留下这一行动节点，前期兵力、粮道和地方协力的次序未逐项保存，不能倒推为单一动因。 「宋金之战：清军占领锦州」之后，朝廷与官员、军队与卫所、边地政权与社群在边地接触区的行动空间有所改写；这种变化须由随后文书和地方材料追踪，不能把初报数字直接当作最终后果。 材料：《崇祯长编》崇祯十五年条及《明史·思宗本纪》；《明史·兵志》及相关列传；边镇、地方志或对方材料。限度：关于「宋金之战：清军占领锦州」，《崇祯长编》崇祯十五年条及《明史·思宗本纪》；《明史·兵志》及相关列传；边镇、地方志或对方材料所能确证的是条目所示的时点和行动；边地接触区内朝廷与官员、军队与卫所、边地政权与社群的规模、意图及社会后果需要异源材料才能判断。
  \item[\eventanchor{ML-1642-002}{1642（崇祯十五年）：浙江遭遇洪涝灾害}] 明与各地社会、边疆政权；朝廷与官员；跨区域行政、资源或军政网络；不假定同步执行。自然风险与地方报告促成朝廷或地方的应对记录。 赈济、迁徙和损失的实际范围仍须以受灾地区材料分别核验。 材料：《崇祯长编》崇祯十五年条及《明史·思宗本纪》；《明史·五行志》；受灾府县地方志与奏疏。限度：以《崇祯长编》崇祯十五年条及《明史·思宗本纪》；《明史·五行志》；受灾府县地方志与奏疏核对「浙江遭遇洪涝灾害」时，只能确认其在1642（崇祯十五年）的叙述位置；朝廷与官员在跨区域行政、资源或军政网络的具体经验仍需由可定位的异类材料补证。
  \item[\eventanchor{ML-1642-003}{1642（崇祯十五年）：复合金属大炮产生于明代。}] 明与各地社会、边疆政权；朝廷与官员；跨区域行政、资源或军政网络；不假定同步执行。制度、教育或知识传播的需要推动了相关行动或文本形成。 它提供制度和传播史的锚点，不能据此推断社会整体接受。 材料：《崇祯长编》崇祯十五年条及《明史·思宗本纪》；《明史·选举志》或《艺文志》；文集、碑刻或书目材料。限度：「复合金属大炮产生于明代」可由《崇祯长编》崇祯十五年条及《明史·思宗本纪》；《明史·选举志》或《艺文志》；文集、碑刻或书目材料定位到1642（崇祯十五年），但这些叙述未完整保存跨区域行政、资源或军政网络的执行与受影响者经验，不能把沉默写成事实。
  \item[\eventanchor{ML-1642-004}{1642（崇祯十五年）：李自成的叛军成功地用大炮在明朝的防御工事上造成了两丈的缺口。}] 明与地方反抗、流动作战集团；朝廷与官员、军队与卫所、地方社会与精英、流动与反抗集团、财政与商业行动者；地方或省域；未具城镇者不补造路线。制度、教育或知识传播的需要推动了相关行动或文本形成。 它提供制度和传播史的锚点，不能据此推断社会整体接受。 材料：《崇祯长编》崇祯十五年条及《明史·思宗本纪》；《明史·选举志》或《艺文志》；文集、碑刻或书目材料。限度：「李自成的叛军成功地用大炮在明朝的防御工事上造成了两丈的缺口」在1642（崇祯十五年）的出现可回查《崇祯长编》崇祯十五年条及《明史·思宗本纪》；《明史·选举志》或《艺文志》；文集、碑刻或书目材料；它们不等于地方或省域的全景记录，朝廷与官员、军队与卫所、地方社会与精英、流动与反抗集团、财政与商业行动者未被记下的选择与损失不作推定。
  \item[\eventanchor{ML-1642-005}{1642（崇祯十五年）：辽饷、剿饷、练饷及地方征解的本年文书显示财政负担，定额、征收与实际流向不得混为一谈。}] 明与各地社会、边疆政权；朝廷与官员、军队与卫所、地方社会与精英、财政与商业行动者；地方或省域；未具城镇者不补造路线。财政征解、交通工程或货币支付的压力使相关措施进入政务。 其法定安排不等于地方执行，后续须以地域材料追踪负担与成效。 材料：《崇祯长编》崇祯十五年条及《明史·思宗本纪》；《明会典》相关条目；地方志、黄册/契约/河工材料。限度：桥接条目只标示可回查的年度问题与材料链；实录有中央选择性，崇祯期尤其需要奏疏、地方、朝鲜及满文材料交叉。
  \item[\eventanchor{ML-1642-006}{1642（崇祯十五年）：朝鲜、辽东和海上关系的本年材料提供外部观察位置，明、后金（清）与地方势力的称谓需具体化。}] 明与各地社会、边疆政权；朝廷与官员、边地政权与社群、地方社会与精英；账本可辨空间；地点细节仍须回到所列材料。朝贡名义、边境安全与港口贸易的利益使交涉持续发生。 它为后续册封、互市或海防安排留下文书与跨政权比较线索。 材料：《崇祯长编》崇祯十五年条及《明史·思宗本纪》；《明史·外国传》；《朝鲜王朝实录》、琉球《历代宝案》或海外档案。限度：桥接条目只标示可回查的年度问题与材料链；实录有中央选择性，崇祯期尤其需要奏疏、地方、朝鲜及满文材料交叉。
  \item[\eventanchor{ML-1642-007}{1642（崇祯十五年）：旱灾、蝗灾、疫病、河患与赈济的本年材料连接环境和社会压力，死亡与流离数字须谨慎处理。}] 明与各地社会、边疆政权；朝廷与官员；跨区域行政、资源或军政网络；不假定同步执行。自然风险与地方报告促成朝廷或地方的应对记录。 赈济、迁徙和损失的实际范围仍须以受灾地区材料分别核验。 材料：《崇祯长编》崇祯十五年条及《明史·思宗本纪》；《明史·五行志》；受灾府县地方志与奏疏。限度：桥接条目只标示可回查的年度问题与材料链；实录有中央选择性，崇祯期尤其需要奏疏、地方、朝鲜及满文材料交叉。
  \item[\eventanchor{ML-1643-001}{1643（崇祯十六年）一月：李自成攻克襄阳}] 明与地方反抗、流动作战集团；朝廷与官员、军队与卫所、地方社会与精英、流动与反抗集团；地方或省域；未具城镇者不补造路线。继承、官僚协调或皇权运作的压力使问题进入朝廷程序。 它影响后续人事与政策议程；动机和派系标签不能由单一叙事裁定。 材料：《崇祯长编》崇祯十六年条及《明史·思宗本纪》；《明史·本纪》及相关列传；诏令、奏疏或会典版本。限度：《崇祯长编》崇祯十六年条及《明史·思宗本纪》；《明史·本纪》及相关列传；诏令、奏疏或会典版本记录「李自成攻克襄阳」的方式限定了可说范围：可追踪1643（崇祯十六年）一月的行动节点，不能凭此还原地方或省域中朝廷与官员、军队与卫所、地方社会与精英、流动与反抗集团的每一处部署。
  \item[\eventanchor{ML-1643-002}{1643（崇祯十六年）十一月：李自成攻克西安}] 明与地方反抗、流动作战集团；朝廷与官员、军队与卫所、地方社会与精英、流动与反抗集团；地方或省域；未具城镇者不补造路线。继承、官僚协调或皇权运作的压力使问题进入朝廷程序。 它影响后续人事与政策议程；动机和派系标签不能由单一叙事裁定。 材料：《崇祯长编》崇祯十六年条及《明史·思宗本纪》；《明史·本纪》及相关列传；诏令、奏疏或会典版本。限度：《崇祯长编》崇祯十六年条及《明史·思宗本纪》；《明史·本纪》及相关列传；诏令、奏疏或会典版本为「李自成攻克西安」留下编年入口，却未必保留地方或省域中朝廷与官员、军队与卫所、地方社会与精英、流动与反抗集团的完整视角；本条因此不据之裁定人数、责任或动机。
  \item[\eventanchor{ML-1643-003}{1643（崇祯十六年）：张献忠在湖广宣告奚朝}] 明与地方反抗、流动作战集团；朝廷与官员、军队与卫所、地方社会与精英、流动与反抗集团；地方或省域；未具城镇者不补造路线。继承、官僚协调或皇权运作的压力使问题进入朝廷程序。 它影响后续人事与政策议程；动机和派系标签不能由单一叙事裁定。 材料：《崇祯长编》崇祯十六年条及《明史·思宗本纪》；《明史·本纪》及相关列传；诏令、奏疏或会典版本。限度：「张献忠在湖广宣告奚朝」借《崇祯长编》崇祯十六年条及《明史·思宗本纪》；《明史·本纪》及相关列传；诏令、奏疏或会典版本获得可检索的时间坐标；地方或省域里朝廷与官员、军队与卫所、地方社会与精英、流动与反抗集团的路线、人数与后果若无其他材料相证，均保留为不可见的部分。
  \item[\eventanchor{ML-1643-004}{1643（崇祯十六年）：辽东防务、后金（清）军事行动和流动作战集团的本年记录标记多战场互动，军数和胜败须保留分歧。（材料核查重点：诏令或奏议的形成与发受链）}] 明与各地社会、边疆政权；朝廷与官员、军队与卫所、边地政权与社群、流动与反抗集团；边地接触区；不把族名或部名当固定疆界。「辽东防务、后金（清）军事行动和流动作战集团的本年记录标记多战场互动，军数和胜败须保留分歧」把朝廷与官员、军队与卫所、边地政权与社群、流动与反抗集团带到边地接触区的具体关口；现存年条只留下这一行动节点，前期兵力、粮道和地方协力的次序未逐项保存，不能倒推为单一动因。 此后，朝廷与官员、军队与卫所、边地政权与社群、流动与反抗集团在边地接触区面对的并非原来的行动条件；「辽东防务、后金（清）军事行动和流动作战集团的本年记录标记多战场互动，军数和胜败须保留分歧」留下需要继续处理的秩序与供给问题，损失和胜负仍应按异源材料分别核对。 材料：《崇祯长编》崇祯十六年条及《明史·思宗本纪》；《明史·兵志》及相关列传；边镇、地方志或对方材料。限度：桥接条目只标示可回查的年度问题与材料链；实录有中央选择性，崇祯期尤其需要奏疏、地方、朝鲜及满文材料交叉。；本条特别核对诏令或奏议的形成与发受链。
  \item[\eventanchor{ML-1643-006}{1643（崇祯十六年）：地方结社、宗教、出版或士人文集的本年线索提供战乱中的社会观察，幸存者记忆不能概括全国。}] 明与各地社会、边疆政权；朝廷与官员、军队与卫所、地方社会与精英、宗教与文化行动者；地方或省域；未具城镇者不补造路线。制度、教育或知识传播的需要推动了相关行动或文本形成。 它提供制度和传播史的锚点，不能据此推断社会整体接受。 材料：《崇祯长编》崇祯十六年条及《明史·思宗本纪》；《明史·选举志》或《艺文志》；文集、碑刻或书目材料。限度：桥接条目只标示可回查的年度问题与材料链；实录有中央选择性，崇祯期尤其需要奏疏、地方、朝鲜及满文材料交叉。
  \item[\eventanchor{ML-1643-007}{1643（崇祯十六年）：天启末至崇祯朝的任免、奏疏和廷议构成本年中枢危机线索；崇祯无实录，必须以多类材料互校。}] 明与各地社会、边疆政权；朝廷与官员；跨区域行政、资源或军政网络；不假定同步执行。继承、官僚协调或皇权运作的压力使问题进入朝廷程序。 它影响后续人事与政策议程；动机和派系标签不能由单一叙事裁定。 材料：《崇祯长编》崇祯十六年条及《明史·思宗本纪》；《明史·本纪》及相关列传；诏令、奏疏或会典版本。限度：桥接条目只标示可回查的年度问题与材料链；实录有中央选择性，崇祯期尤其需要奏疏、地方、朝鲜及满文材料交叉。
  \item[\eventanchor{ML-1644-001}{1644（崇祯十七年）二月8日：李自成在西安宣告大顺}] 明与地方反抗、流动作战集团；朝廷与官员、军队与卫所、地方社会与精英、流动与反抗集团；地方或省域；未具城镇者不补造路线。继承、官僚协调或皇权运作的压力使问题进入朝廷程序。 它影响后续人事与政策议程；动机和派系标签不能由单一叙事裁定。 材料：《崇祯长编》崇祯十七年条及《明史·思宗本纪》；《明史·本纪》及相关列传；诏令、奏疏或会典版本。限度：《崇祯长编》崇祯十七年条及《明史·思宗本纪》；《明史·本纪》及相关列传；诏令、奏疏或会典版本为「李自成在西安宣告大顺」留下编年入口，却未必保留地方或省域中朝廷与官员、军队与卫所、地方社会与精英、流动与反抗集团的完整视角；本条因此不据之裁定人数、责任或动机。
  \item[\eventanchor{ML-1644-002}{1644（崇祯十七年）四月25日：李自成占领北京，崇祯皇帝上吊自杀}] 明与地方反抗、流动作战集团；朝廷与官员、军队与卫所、地方社会与精英、流动与反抗集团；地方或省域；未具城镇者不补造路线。继承、官僚协调或皇权运作的压力使问题进入朝廷程序。 它影响后续人事与政策议程；动机和派系标签不能由单一叙事裁定。 材料：《崇祯长编》崇祯十七年条及《明史·思宗本纪》；《明史·本纪》及相关列传；诏令、奏疏或会典版本。限度：就「李自成占领北京，崇祯皇帝上吊自杀」而论，《崇祯长编》崇祯十七年条及《明史·思宗本纪》；《明史·本纪》及相关列传；诏令、奏疏或会典版本提供的是1644（崇祯十七年）四月25日的记录线索；朝廷与官员、军队与卫所、地方社会与精英、流动与反抗集团在地方或省域的实际行动细节不能由这一材料链独自补足。
  \item[\eventanchor{ML-1644-003}{1644（崇祯十七年）三月至四月：京畿防务、流动作战与军饷决策的文书显示崇祯十七年初的中枢危机；各路兵力和命令传递需要多源校验。}] 明与地方反抗、流动作战集团；朝廷与官员、军队与卫所、地方社会与精英、流动与反抗集团；地方或省域；未具城镇者不补造路线。辽东战事、财政征解和流动作战集团的压力同时聚集于京畿。 其记录为考察李自成入京前的防务与行政失灵提供可交叉核验的材料链。 材料：《崇祯长编》崇祯十七年年初条；《明季北略》；《明史·思宗本纪》及京畿地方材料。限度：崇祯无实录；近时见闻、遗民记忆与胜者叙事均有立场，必须将日期、军数和命令传达分别核验。
  \item[\eventanchor{ML-1644-004}{1644（崇祯十七年）：朝鲜、辽东和海上关系的本年材料提供外部观察位置，明、后金（清）与地方势力的称谓需具体化。（材料核查重点：诏令或奏议的形成与发受链）}] 明与各地社会、边疆政权；朝廷与官员、边地政权与社群、地方社会与精英；账本可辨空间；地点细节仍须回到所列材料。朝贡名义、边境安全与港口贸易的利益使交涉持续发生。 它为后续册封、互市或海防安排留下文书与跨政权比较线索。 材料：《崇祯长编》崇祯十七年条及《明史·思宗本纪》；《明史·外国传》；《朝鲜王朝实录》、琉球《历代宝案》或海外档案。限度：桥接条目只标示可回查的年度问题与材料链；实录有中央选择性，崇祯期尤其需要奏疏、地方、朝鲜及满文材料交叉。；本条特别核对诏令或奏议的形成与发受链。
  \item[\eventanchor{ML-1644-005}{1644（崇祯十七年）：朝鲜、辽东和海上关系的本年材料提供外部观察位置，明、后金（清）与地方势力的称谓需具体化。（材料核查重点：报称信息的来源、抄传与行政分类）}] 明与各地社会、边疆政权；朝廷与官员、边地政权与社群、地方社会与精英；账本可辨空间；地点细节仍须回到所列材料。朝贡名义、边境安全与港口贸易的利益使交涉持续发生。 它为后续册封、互市或海防安排留下文书与跨政权比较线索。 材料：《崇祯长编》崇祯十七年条及《明史·思宗本纪》；《明史·外国传》；《朝鲜王朝实录》、琉球《历代宝案》或海外档案。限度：桥接条目只标示可回查的年度问题与材料链；实录有中央选择性，崇祯期尤其需要奏疏、地方、朝鲜及满文材料交叉。；本条特别核对报称信息的来源、抄传与行政分类。
  \item[\eventanchor{ML-1644-006}{1644（崇祯十七年）：旱灾、蝗灾、疫病、河患与赈济的本年材料连接环境和社会压力，死亡与流离数字须谨慎处理。}] 明与各地社会、边疆政权；朝廷与官员；跨区域行政、资源或军政网络；不假定同步执行。自然风险与地方报告促成朝廷或地方的应对记录。 赈济、迁徙和损失的实际范围仍须以受灾地区材料分别核验。 材料：《崇祯长编》崇祯十七年条及《明史·思宗本纪》；《明史·五行志》；受灾府县地方志与奏疏。限度：桥接条目只标示可回查的年度问题与材料链；实录有中央选择性，崇祯期尤其需要奏疏、地方、朝鲜及满文材料交叉。
  \item[\eventanchor{ML-1644-007}{1644（崇祯十七年）：辽饷、剿饷、练饷及地方征解的本年文书显示财政负担，定额、征收与实际流向不得混为一谈。}] 明与各地社会、边疆政权；朝廷与官员、军队与卫所、地方社会与精英、财政与商业行动者；地方或省域；未具城镇者不补造路线。财政征解、交通工程或货币支付的压力使相关措施进入政务。 其法定安排不等于地方执行，后续须以地域材料追踪负担与成效。 材料：《崇祯长编》崇祯十七年条及《明史·思宗本纪》；《明会典》相关条目；地方志、黄册/契约/河工材料。限度：桥接条目只标示可回查的年度问题与材料链；实录有中央选择性，崇祯期尤其需要奏疏、地方、朝鲜及满文材料交叉。
\end{description}
